\documentclass[11pt,a4paper]{article}
% Standalone research plan, 24 September 2026 (version 7.1, full-time).
% Compile with latexmk -pdf, or run pdflatex twice.
% Model assumptions, derivations, completed checks and bibliography are included below.
\usepackage[utf8]{inputenc}
\usepackage[T1]{fontenc}
\input{glyphtounicode}\pdfgentounicode=1
\IfFileExists{lmodern.sty}{\usepackage{lmodern}\usepackage{microtype}}{\usepackage[expansion=false]{microtype}}
\usepackage[margin=2.5cm]{geometry}
\usepackage{amsmath,amssymb,booktabs,array,longtable,enumitem,xcolor,framed}
\usepackage[colorlinks=true,linkcolor=blue!55!black,citecolor=blue!55!black,urlcolor=blue!55!black]{hyperref}
\usepackage{bookmark,needspace}
\newcolumntype{P}[1]{>{\raggedright\arraybackslash}p{#1}}
\newcommand{\E}{\mathbb E}
\newcommand{\Var}{\operatorname{Var}}
\newcommand{\eps}{\varepsilon}
\newcommand{\dd}{\mathrm d}
\newcommand{\KL}{\operatorname{KL}}
\newcommand{\DL}{\mathrm{DL}}
\newcommand{\EDL}{\mathrm{EDL}}
\newcommand{\IIA}{\mathrm{IIA}}
\newcommand{\unv}{\textsuperscript{\dag}}
\definecolor{cknown}{RGB}{20,100,40}
\definecolor{ccalc}{RGB}{25,70,150}
\definecolor{cnum}{RGB}{0,105,115}
\definecolor{chyp}{RGB}{150,30,40}
\definecolor{ctest}{RGB}{175,85,0}
\definecolor{cassess}{RGB}{95,70,125}
\newcommand{\evl}[2]{{\small\bfseries\color{#1}[#2]}}
\newcommand{\LKnown}{\evl{cknown}{Known result}}
\newcommand{\LCalc}{\evl{ccalc}{Internal calculation}}
\newcommand{\LNum}{\evl{cnum}{Numerical check}}
\newcommand{\LHyp}{\evl{chyp}{Hypothesis}}
\newcommand{\LTest}{\evl{ctest}{Proposed test}}
\newcommand{\LAssess}{\evl{cassess}{Assessment}}
\newenvironment{hypbox}[1]{%
  \def\FrameCommand{{\color{chyp}\vrule width 1.6pt}\hspace{8pt}}%
  \MakeFramed{\advance\hsize-\width\FrameRestore}%
  \noindent\textbf{#1}\par\smallskip\noindent\ignorespaces}{\endMakeFramed}
\newenvironment{keybox}[1]{%
  \def\FrameCommand{{\color{ccalc}\vrule width 1.6pt}\hspace{8pt}}%
  \MakeFramed{\advance\hsize-\width\FrameRestore}%
  \noindent\textbf{#1}\par\smallskip\noindent\ignorespaces}{\endMakeFramed}
\setlist{itemsep=3pt,topsep=5pt}
\setlength{\emergencystretch}{2em}
\setcounter{tocdepth}{1}
\title{\textbf{Bits of Evidence}\\[0.4em]
\large When description lengths expose manufactured explanations,\\
and what singular learning theory says about the codes\\[0.3em]
\normalsize An 18-month research plan, version 7.1 (full-time)}
\author{}
\date{24 September 2026}
\begin{document}
\maketitle
\begin{abstract}
\noindent
Interpretability and safety research make evidential claims about models,
and the tools used to support those claims can manufacture them. An
expressive alignment map can make an untrained network ``implement'' an
algorithm; an elicitation procedure can ``reveal'' a capability it in fact
teaches. This plan asks when information-theoretic measurements can tell
genuine structure from structure supplied by the tool, and what singular
learning theory (SLT) says about the codes involved. Log loss is a code
length, sequential Bayesian prediction gives the mixture code whose
redundancy SLT describes, and likelihood ratios support tests once their
null distributions and sequential construction are specified.

The flagship, Project~A, proceeds in two stages. Stage~1 measures
\emph{alignment sample efficiency}: how many bits a specified alignment
learner needs to make a network reproduce a fixed algorithm's
counterfactual targets. Comparisons are made across trained and untrained
networks, loci and map families on the same targets. The score is
validated against distribution-shifted held-out accuracy and against
causal-abstraction error with interventions on unmapped variables, on
ground-truth models and on false explanations built to pass existing
tests. Stage~2 is the ambitious goal: explanations as codes for the
network's own interventional behaviour, in which competing causal
abstractions, data-driven interfaces and capacity-matched black boxes
predict the same network responses. It is positioned against recent work
on prequential causal description length and applied to refusal and to
deception probes.

A theory strand establishes what prequential scores measure. Block codes
multiply logarithmic coefficients by a computable batching factor;
plug-in and gradient-trained predictors need not recover the learning
coefficient; prior-dependent constants can dominate at attainable sample
sizes; Bayes factors are e-values only against specified priors; and
guaranteed tests against composite control families require certified
likelihood bounds. SLT hypotheses are tested where the learning
coefficient is known.

Project~B, gated on validation, asks whether one independently measured
geometric statistic---a local free-energy depth, first validated on small
models with exactly computable basin masses---predicts whether a
capability is recoverable under a stated adaptation protocol better than
excess description length and calibrated uplift. Stretch goals include
organisms with seven to eight billion parameters, susceptibility-based
predictions of unlock cost and rare-output estimation.

The plan opens with a six-week decision experiment and a pre-registered
gate. It keeps its bold hypotheses, labels them as such, and pairs every gate with a stop rule. The document separates known results,
internal calculations, completed numerical checks, hypotheses and proposed
tests.
\end{abstract}
\newpage
{\normalsize\tableofcontents}
\newpage

\section{Purpose, constraints and what changed}
\label{sec:purpose}

\subsection{The objective}

Eighteen months remain on the current contract. The work should position
the author for an AI-safety or interpretability role afterwards, in
industry or academia. Applications will be written around months 12--15,
so the visible record---preprints, posts and code---must exist by month
12. Industry teams look for competent end-to-end experiments on real
language models, at least one result the field discusses, and a
distinctive strength. Academic committees look for main-venue papers with
a coherent theme and theoretical depth. The distinctive strength here is
calibrated statistical measurement: code lengths, Bayesian model
comparison, sampler diagnostics and honest uncertainty.

\paragraph{Working assumptions.} The author works on this programme full
time: about twelve person-months before month 12 and eighteen in total.
Compute is available and does not constrain the design. The binding
constraints are the decision gates, engineering at the scale of billions
of parameters, and the calendar of review cycles. Month 0 is October 2026.

\subsection{What an external review changed}

In September 2026 an external review of version 6.4 found that the move to
description lengths puts SLT on firmer ground. It also found that the
flagship's central score did not measure what the surrounding prose
claimed, that close prior work was missing, that several bridges from
operational measurements to SLT or to statistical guarantees were wrong,
and that Project~B combined too many uncertain ideas. The review's factual
claims were checked. The prior work exists and is close \cite{ma2026};
excess description length subtracts the final model's test loss, not its
training loss \cite{donoway2026}; and Sutter et al.\ already evaluate on
held-out data \cite{sutter2025}. Table~\ref{tab:review} records how each
point is treated. No point was rejected; two were modified.

\begin{longtable}{@{}P{5.0cm}P{1.9cm}P{7.8cm}@{}}
\caption{Treatment of the external review.}\label{tab:review}\\
\toprule
Review point & Status & Treatment in version 7 \\
\midrule
\endfirsthead
\toprule
Review point & Status & Treatment in version 7 \\
\midrule
\endhead
The flagship score codes the algorithm's labels with the network, so it
measures alignment sample efficiency, not bits of network behaviour
explained & Adopted & Stage~1 is named and scoped accordingly; coding the
network's own behaviour becomes Stage~2 (Sections~\ref{sec:stage1}
and~\ref{sec:a1a2}). \\ \addlinespace
Different candidate algorithms generate different target sets, so their
code lengths are not comparable & Adopted & Stage~1 compares networks, loci
and map families on fixed targets only; comparisons between algorithms
move to Stage~2's common responses. \\ \addlinespace
A network-free control that is given the algorithm achieves the minimum
code length & Adopted & Every control's information access is specified;
the label-only baseline has no access to the algorithm. \\ \addlinespace
Close prior work on prequential causal description length & Adopted &
Positioning rewritten; Stage~2 is designed as a comparison with, and
extension of, \cite{ma2026}. \\ \addlinespace
Held-out evaluation is already part of \cite{sutter2025} & Adopted,
extended & Distribution-shifted held-out accuracy, which already breaks
some manufactured alignments \cite{sutter2025}, becomes a primary rival. \\
\addlinespace
Smaller alignment maps do not imply smaller learning coefficients;
elicitation need not mean small $\lambda$ & Adopted & Both are stated as
hypotheses; prior-dependent constant terms are part of the analysis. \\
\addlinespace
Perfect IIA does not imply equal log loss & Adopted & Faithfulness is
measured by log loss with a calibration check, and the regimes are defined
by log loss. \\ \addlinespace
Doubling blocks change the logarithmic coefficient & Adopted, extended & A
general batching factor is derived and checked
(Appendix~\ref{app:batch}); SLT comparisons use per-sample Bayes codes or
corrected block codes. \\ \addlinespace
EDL subtracts the final test loss & Adopted & Definition and decomposition
corrected (Appendix~\ref{app:edl}). \\ \addlinespace
E-values need conditional validity; universal inference needs the
maximized null likelihood & Adopted & Guaranteed tests are restricted to
control families with certified likelihood bounds, and sequential
constructions use fresh data (Appendix~\ref{app:evidence}). \\ \addlinespace
Complexity penalties cannot replace interventions on unmapped variables &
Adopted & The intervention protocol includes unmapped variables. \\
\addlinespace
Project~B combines too many uncertain ideas; basin masses are hard to
measure; ``suppressed'' needs an operational definition; adding data and
replacing data differ & Adopted & Project~B is narrowed to one geometric
statistic, validated first on small models with exact masses. Conclusions
concern recoverability under a stated protocol; the addition formula is
corrected and the flip is computed exactly by reweighting
(Appendix~\ref{app:flip}). \\ \addlinespace
Replace the three-paper commitment by a six-week decision experiment;
allow equivalent valid abstractions & Adopted & Weeks 1--6 end in the
review's gate criterion (Section~\ref{sec:gatesA}). \\ \addlinespace
Remove numerical publication probabilities & Modified & The plan uses
scenarios and readiness indicators (Section~\ref{sec:outlook});
subjective probabilities are kept outside the plan. \\ \addlinespace
Do not promise separate papers in advance & Modified & Papers are
conditional on gates, but submission targets remain, because the record
matters for the objective. \\
\bottomrule
\end{longtable}

\subsection{Why the plan remains ambitious}

Narrowing the flagship's first claim does not lower its target. The goal
is a validated, theoretically understood method for telling when an
interpretability tool supplies the apparent explanation itself. Current
practice cannot answer that question, and any safety argument that rests
on interpretability depends on it. Stage~2 aims further, at a general way
to score explanations as codes for a network's interventional behaviour.
The theory strand aims at a precise account of what prequential scores
measure, with SLT describing the Bayesian part. Project~B aims at an audit
statistic whose validity is established rather than assumed. Each of these
can fail. The plan accepts that risk, states its bold hypotheses as
hypotheses, and pairs every gate with a stop rule, so that a failing direction is detected early, stopped and shared rather than extended.

\Needspace{7\baselineskip}
\paragraph{How to read this plan.} Section~\ref{sec:primer} explains the
ideas needed later. Section~\ref{sec:projA} is the flagship with its theory
strand, Section~\ref{sec:projB} the gated second project,
Section~\ref{sec:p0} a short note on training dynamics, and
Sections~\ref{sec:tools}--\ref{sec:risks} cover tooling, schedule, outlook
and risks. Appendices collect derivations and completed checks.

\Needspace{7\baselineskip}
\paragraph{What the evidence labels mean.}
\begin{description}[style=nextline,leftmargin=1.5em]
\item[\LKnown] Established in the literature; the source is cited where
used.
\item[\LCalc] Worked out in this document under stated assumptions. The
label implies neither novelty nor a completed proof.
\item[\LNum] A completed computation, archived with code, that tests a
displayed formula. It concerns the stated mathematical model, not observed
training of large networks.
\item[\LHyp] A claim the plan expects to be true and designs tests for.
\item[\LTest] An experiment or theorem still to be completed.
\item[\LAssess] A subjective judgement.
\end{description}
\noindent A dagger\unv{} on a reference marks bibliographic details that
were not independently verified when this version was written.

\section{Primer: bits, codes and singular learning theory}
\label{sec:primer}

\subsection{Log loss is a code length}

\LKnown{} A predictive distribution $q$ that assigns probability $q(y)$
to an outcome can be turned into a prefix code whose codeword for $y$
has about $-\log_2 q(y)$ bits (Kraft's inequality), and no code does
better on average than the entropy \cite{grunwald2007}. The cumulative
log loss of a sequence of predictions is therefore, up to rounding, the
length of a code for the observed labels. Formulas below use natural
logarithms (nats); prose reports bits, with $1$ nat $\approx 1.44$ bits.

\paragraph{Example.} Encoding 1000 fair-coin outcomes costs 1000 bits;
nothing can be saved. If the coin in fact lands heads with probability
0.9, a code that knows this needs about $1000\,h(0.9)\approx469$ bits,
where $h$ is the binary entropy. A learner that must first discover the
bias pays those 469 bits plus a small learning cost. The learning cost
is where model complexity appears.

\subsection{Three codes for the same data}

\LKnown{} Let $D_n=\{(x_i,y_i)\}_{i\le n}$ be labelled data and
$\{p(y\mid x,w)\}$ a model with prior $\varphi$.
\begin{enumerate}[leftmargin=*]
\item A \emph{two-part code} first describes a hypothesis, then the labels
given the hypothesis.
\item The \emph{Bayes mixture code} uses the marginal likelihood; its length
is the Bayes free energy
\begin{equation}
F_n=-\log\int\prod_{i=1}^n p(y_i\mid x_i,w)\,\varphi(w)\,\dd w .
\label{eq:freeenergy}
\end{equation}
\item A \emph{prequential code} \cite{dawid1984} predicts each label with a
model fitted to the previous labels only:
$L^A(n)=\sum_{i=1}^n-\log p^A_{i-1}(y_i\mid x_i)$ for a learning
algorithm $A$.
\end{enumerate}
By the chain rule, the prequential code built from the Bayes predictive
distribution is exactly the mixture code. Prequential codes built from
SGD-trained networks are valid codes too \cite{blierollivier2018}, but
they are not Bayes codes; the difference is analysed in
Section~\ref{sec:projB}. In practice the data are coded in blocks: fit
on the first $k$ blocks, code block $k+1$, and repeat
\cite{voita2020,blierollivier2018}.

\subsection{How complexity enters: regular versus singular models}

\LKnown{} For a regular model with $d$ identifiable parameters and
positive-definite Fisher information,
$F_n=nL_n(w_0)+\tfrac d2\log n+O_p(1)$, where $L_n$ is the average
negative log likelihood and $w_0$ the optimal parameter
\cite{rissanen1986,clarkebarron1990}. Neural networks are singular:
many parameters give the same function and the Fisher matrix is
degenerate. Watanabe proved
\begin{equation}
F_n=nL_n(w_0)+\lambda\log n-(m-1)\log\log n+O_p(1),
\label{eq:watanabe}
\end{equation}
where $\lambda\le d/2$ is the learning coefficient (real log canonical
threshold) and $m$ its multiplicity \cite{watanabe2009}. Equivalently,
the expected excess log loss of the Bayes predictive at sample size $n$
is $\lambda/n$. Two-part codes with the same $\lambda\log n$ redundancy
exist for singular models, and the local learning coefficient tracks
the compressibility of trained language models \cite{urdshals2025}.
Local versions of $\lambda$ at a particular solution are estimated by stochastic-gradient Langevin dynamics \cite{lau2025llc}. The constant term depends on the prior: two regular models of the same dimension share $\lambda=d/2$ even when their priors make one target far easier to learn, so finite-sample code lengths can differ in terms that $\lambda$ does not see.

\paragraph{Worked example.} \LCalc{} Take $y=abx+\text{noise}$ with two
parameters $(a,b)$. Counting parameters suggests a complexity term
$\tfrac d2\log n=\log n$. But the parameters with the correct product
form a curve when the true coefficient $c^*=ab$ is nonzero, and the
union of two axes when $c^*=0$. Both cases have $\lambda=1/2$, with
$m=1$ and $m=2$ respectively. At $n=10^4$ the complexity term is
$\tfrac12\ln10^4\approx4.6$ nats ($6.6$ bits) instead of $9.2$ nats
($13.3$ bits). Parameter counting doubles the cost here. For the
overparameterized families used as alignment maps and adapters the gap
can be much larger; reduced-rank regression is a solved example
\cite{aoyagiwatanabe2005}.

\subsection{When a family cannot compress: minimax regret}

\LKnown{} For a family $\mathcal H$ of predictors and fixed inputs, the
smallest achievable worst-case regret of a code relative to the best
member in hindsight is
$\log\sum_{y^n}\sup_{h\in\mathcal H}p_h(y^n\mid x^n)$, the normalizer of
the normalized maximum likelihood code \cite{shtarkov1987,grunwald2007}.
If a family of deterministic predictors realizes every labelling of $n$
items over an alphabet $\mathcal Y$, the sum has $|\mathcal Y|^n$ unit
terms and the regret is $n\log|\mathcal Y|$. A family that can fit
anything cannot compress anything. This is the information-theoretic
content of the dilemma addressed in Section~\ref{sec:projA}.

\subsection{Local free energy and posterior mass}

\LKnown{} For a region $U$ of parameter space, the local free energy
\begin{equation*}
F_U(n)=-\log\int_U e^{-nL_n(w)}\varphi(w)\,\dd w
\end{equation*}
has the expansion
\eqref{eq:watanabe} with local values $(\lambda_U,m_U)$ plus a constant
$c_U$ \cite{watanabe2009}. Posterior odds between regions are
$\exp[-(F_{U_1}-F_{U_2})]$. Tempering at inverse temperature $\beta$
replaces $n$ by $n\beta$. Stochastic-gradient MCMC samples localized,
tempered posteriors; the estimator $\hat\lambda=n\beta\,(\E_\beta
L_n-L_n(w^*))$ is standard \cite{lau2025llc}. Its reliability depends
on sampler settings and on whether $w^*$ is a minimum
\cite{hitchcock2025,cai2026}. Posterior-covariance statistics of this
kind have been estimated for networks with billions of parameters
\cite{kreer2026}.

\subsection{Excess description length}

\LKnown{} Excess description length (EDL) is the prequential code length of the fine-tuning labels under online training, minus $n$ times the final model's expected loss on fresh data, its test loss. Subtracting test loss rather than training loss means that memorizing random labels earns no credit. EDL is nonnegative in expectation under conditions on the learning procedure, converges to surplus description length, and is defined relative to a specific training algorithm \cite{donoway2026}. An
empirical companion reports that eliciting a latent capability needs
orders of magnitude less information than teaching a new one
\cite{bitsthatcount2026}. A related study fits a logistic curve relating
the number of trained parameters to recovery of a capability gap
\cite{elicitation2025}.

\subsection{E-values: evidence with error guarantees}
\label{sec:evalues}

\LKnown{} An e-value for a null hypothesis is a nonnegative statistic $E$
whose expectation is at most one under every distribution in the null. By
Markov's inequality, rejecting when $E\ge1/\alpha$ has type-I error at most
$\alpha$. Unlike p-values, e-values keep this guarantee under optional
continuation---deciding to collect more data after seeing earlier
results---and can be combined across studies
\cite{grunwald2024safe,ramdas2023}. A running likelihood ratio that is a
nonnegative martingale under the null gives anytime-valid monitoring
through Ville's inequality. Two constructions matter here. A Bayes factor
$q_1(y^n)/q_0(y^n)$ between two mixture codes is an e-value against the
simple null that the data follow $q_0$, and growth-rate-optimal e-values for
composite nulls also take the form of Bayes factors with special priors
\cite{grunwald2024safe}. Universal inference turns any prequential likelihood, or one fitted on separate data, into an e-value against a composite null by dividing it by the null's maximized likelihood \cite{wasserman2020}. The maximum must be computed exactly or bounded from above; an optimizer's best value is only a lower bound and would inflate $E$. Multiplying e-values from successive rounds requires each to be valid conditionally on the past, $\E[E_t\mid\mathcal F_{t-1}]\le1$ under the null. Sequential likelihood ratios satisfy this by construction under the null that defines them; repeated calibrations of the same audited object need not, even with fresh data. In bits, $\log_2E$ is the
evidence: $E=2^{10}$ rejects at level about $10^{-3}$. Universal codes, the backbone of MDL, are natural sources of e-values with small regret; unifying description-length model selection with e-value testing is an open problem.

\subsection{Susceptibilities: how the local posterior responds to data}
\label{sec:suscept}

\LKnown{} Susceptibilities treat a trained network as a Bayesian
statistical-mechanical system \cite{baker2026}. Shift the data distribution
slightly from $q$ towards $q'$ and ask how the posterior expectation of an
observable $\phi$ localized on a component responds. \LCalc{} With
$q_h=(1-h)q+hq'$ and the tempered local posterior proportional to
$\exp[-n\beta L_h(w)]\varphi(w)$,
\begin{equation}
\chi_\phi=\frac{\dd}{\dd h}\E_h[\phi]\Big|_{h=0}
=-n\beta\,\mathrm{Cov}\bigl(\phi,\,L_{q'}-L_q\bigr),
\label{eq:suscept}
\end{equation}
up to the normalization conventions of \cite{baker2026}. This describes a shift at fixed data size; adding data instead changes the formula (Appendix~\ref{app:lr}). The covariance is
estimated from local SGLD samples and splits into per-token contributions.
Response matrices built from many shifts separate functional modules such
as induction and multigram heads \cite{baker2026}, and clusters found in a
14M-parameter model remain coherent in models up to 1.4B parameters
\cite{gordon2026}. Inverting the linear response gives \emph{patterning}:
choosing training data that steers internal structure, including which of
several perfect algorithms a small model learns, by targeting the local
learning coefficient of each solution \cite{wangmurfet2026}; a follow-up
applies the idea to debiasing reward models \cite{patterningpractice2026}.
Informally, \eqref{eq:suscept} says how strongly data from $q'$ would move the posterior's view of $\phi$, that is, how informative such data would be about that component.

\subsection{Causal abstraction in one page}

\LKnown{} A high-level causal model, here an algorithm $A$ with
intermediate variables $V$, is related to a network $N$ through an
\emph{interchange intervention}. Run $N$ on a base input $b$, but
replace the part of an internal representation aligned with $V$ by its
value on a source input $s$. The algorithm predicts the counterfactual
output $A(b;V\leftarrow V(s))$. Interchange-intervention accuracy (IIA)
is the fraction of pairs on which the intervened network agrees with
that prediction \cite{geiger2025ca}. Distributed alignment search learns
the alignment, typically a rotation, by gradient descent \cite{geiger2024das}.

Sutter et al.\ proved that unrestricted alignment maps make any network
implement any algorithm under mild conditions. With nonlinear maps they
reached perfect IIA on the indirect-object-identification task in
randomly initialized language models, and they conclude that no principled way currently exists to trade map complexity against accuracy \cite{sutter2025}. Their maps generalize to unseen test pairs, but when training and test pairs use disjoint sets of names, the alignment learned on an untrained model fails to generalize: distribution-shifted held-out tests already expose some manufactured alignments. They also observed that the map complexity needed for
perfect alignment falls as training proceeds. Causal abstractions need
not be unique \cite{meloux2025ident}. A new benchmark of ten simulated
systems with valid and invalid abstractions evaluated more than thirty
metrics; only causal metrics that also test faithfulness on unmapped
variables discriminated reliably \cite{meloux2026}. Ground-truth
testbeds exist: compiled transformers \cite{tracr2023}, semi-synthetic
transformers with known circuits \cite{interpbench2024}, and a benchmark
whose causal-variable track found distributed alignment search best \cite{mib2025}. Closest to this plan, Ma et al.\ compare candidate interface architectures in language models by prequential causal description length over ``forked futures''---response distributions of operations sampled after a hidden state has formed---and support their conclusions with transplantation, mediation and a blind model-organism test \cite{ma2026}.

\subsection{Hidden capabilities in one page}

\LKnown{} Language models can be prompted or fine-tuned to underperform
strategically (``sandbagging'') \cite{vanderweij2024}.
Password-locked models, which perform well only when a trigger is
present, are standard proxies \cite{greenblatt2024}. Unlearning methods
aim to remove hazardous knowledge \cite{li2024wmdp}, but removed
knowledge is often recoverable \cite{lucki2024}. Injecting noise into
the weights raises the accuracy of sandbagging models while honest
models degrade, which gives a detection signal \cite{tice2025}. In an
adversarial auditing game, however, a blinded team could not reliably
separate sandbagging from benign models. Probes were promising but were
confused by distractor behaviours, and fine-tuning on a single
demonstration elicited every sandbagger but also improved benign models
\cite{taylor2025}. What is missing is a calibrated null: how much
improvement, or how few bits, should we expect from a model that is not
hiding anything?

\paragraph{Rare outputs.} \LKnown{} A related problem is estimating the
probability of a harmful output too rare to observe by sampling.
Importance sampling outperforms extrapolation from activations for small
transformers \cite{wuhilton2025}, and proposal models built by activation
steering estimate rare harmful, hallucinated and sycophantic outputs with
about 5\% of the samples plain Monte Carlo needs \cite{angell2026}.

\section{Project A (flagship): when do description lengths expose manufactured explanations?}
\label{sec:projA}

\subsection{The question}

An interpretability claim is only as good as the analysis tool's inability
to produce it spuriously. Sutter et al.\ showed that expressive alignment
maps make causal abstraction vacuous \cite{sutter2025}. Méloux et al.\
showed that most metrics fail to separate valid from invalid abstractions
unless they test interventions on unmapped variables \cite{meloux2026}. Ma
et al.\ proposed prequential causal description length for choosing between
interface architectures \cite{ma2026}. The flagship asks the question these
results leave open: \emph{under what conditions do description-length
measurements distinguish genuine causal structure from alignments
manufactured by an expressive analysis tool, and what exactly do the
measurements quantify?} A strong answer combines positive results, clear
counterexamples, a validated protocol and a theory of the codes. It matters
for safety because probes, directions and circuits increasingly appear in
safety arguments, and a method that cannot tell when the tool supplied the
explanation cannot certify them.

\subsection{Stage 1: alignment sample efficiency}
\label{sec:stage1}

\LCalc{} Fix an algorithm $A$ with intermediate variables $V$ and a
distribution over base and source inputs $(b,s)$. The targets are the
algorithm's counterfactual outputs $y^*(b,s)=A(b;V\leftarrow V_A(s))$,
passed through randomized response with a known small noise rate. A
\emph{configuration} is a network $N$, a locus $\ell$ and an alignment
learner, that is, a map family $\Phi$ with a prior or training procedure.
Its prequential code length on the fixed targets is
\begin{equation}
\DL(n)=\sum_{i=1}^n-\log N_{\hat\theta_{i-1}}(b_i,s_i)\bigl[y^*_i\bigr],
\label{eq:dlh}
\end{equation}
where $N_\theta(b,s)$ is the output distribution of $N$ on $b$ with the
aligned part of $h_\ell(b)$ replaced by its value on $s$, and
$\hat\theta_{i-1}$ is fitted on the first $i-1$ pairs.

\paragraph{What the score measures.} $\DL(n)$ is the cost, in bits, for a
specified alignment learner to make this network reproduce this
algorithm's counterfactual targets. It does not count bits of the
network's behaviour explained by the algorithm: the labels come from the
algorithm, and the network serves as the predictor. The score becomes
informative through comparisons on the \emph{same} targets. The gain
$G(n)=\DL_{\rm control}(n)-\DL(n)$ compares a configuration with a control
that differs in one respect: an untrained network of the same
architecture, a locus before $V$ can be computed, or a different map
family. Every learner sees the inputs only through the network's
representation at its locus and the network's fixed readout, and none has
access to $A$. A label-only code that sees neither the network nor $A$
normalizes for the difficulty of the targets.

\paragraph{Faithfulness in log loss.} Interchange-intervention accuracy
(IIA) measures agreement of the most probable output. Two configurations
can both reach perfect accuracy while assigning the correct label
probabilities 0.51 and 0.99. The plan therefore measures faithfulness by
log loss, checks calibration, and defines the regimes below by log loss
rather than accuracy.

\begin{keybox}{Decomposition on fixed targets (internal calculation).}
For Bayes codes over families that satisfy the conditions behind
\eqref{eq:watanabe},
\begin{equation}
G(n)=n\,\Delta\ell+\Delta\lambda\log n-\Delta(m-1)\log\log n+R_n+O_p(1),
\label{eq:gain}
\end{equation}
where $\Delta\ell$ is the gap in optimal per-pair log loss, $\Delta\lambda$
and $\Delta m$ are the gaps in learning coefficient and multiplicity at the
respective optima, and $R_n$ is the difference of centred empirical losses.
If both configurations reproduce the law of the randomized-response targets
exactly, they make identical predictions and $\Delta\ell=R_n=0$. The
$O_p(1)$ terms depend on the priors and need not be small at attainable
$n$: two regular families of equal dimension share $\lambda=d/2$ even when
their priors make one target much easier to learn. Derivation:
Appendix~\ref{app:codes}.
\end{keybox}

\paragraph{Three regimes.} \LCalc{} If the control cannot match the
configuration's log loss ($\Delta\ell>0$), the gain grows linearly and the
control can be rejected: the falsification regime. If both reproduce the same conditional target law, the gain grows at most logarithmically: the Occam regime. There, no valid test can reject a
control family that contains the target law, and the evidence is
necessarily relative to the priors or learners involved (Section~\ref{sec:atests}). Equal expected log loss is not enough: two different predictors with equal expected loss leave the fluctuation term $R_n$, which can grow like $\sqrt n$, and calibration alone does not make their conditional predictions equal. If neither fits, the comparison concerns
approximation quality and is read together with the faithfulness gap.
Untrained networks with expressive maps are candidates for the Occam
regime on i.i.d.\ splits; distribution-shifted splits can move them into
the falsification regime \cite{sutter2025}.

\subsection{What prequential scores measure}
\label{sec:atests}

The theory strand makes four statements precise; each has a proof target
below.
\begin{description}[style=nextline,leftmargin=1.5em]
\item[Batching.] \LCalc{} If block boundaries grow geometrically by a
factor $r$ and the learner's expected excess log loss after $m$ examples is
$\kappa/m$, the cumulative excess is $\kappa\frac{r-1}{\ln r}\log n+o(\log n)$, with an $O(1)$ remainder when the learning curve's own remainder is $O(m^{-1-\delta})$ for some $\delta>0$.
Doubling blocks multiply the coefficient by $1/\ln2\approx1.44$, even for
Bayes predictors frozen within blocks; per-sample updating recovers $\kappa\log n$. In singular models with multiplicity above one, the Bayes learning curve carries a term of order $1/(m\log m)$, which accumulates into a $\log\log n$ correction. \LNum{} Exact computations for a Gaussian mean confirm the
factor to three decimals for plug-in and Bayes predictors
(Appendix~\ref{app:batch}).
\item[Learner coefficients.] \LKnown{} For the Bayes predictive,
$\kappa=\lambda$ \cite{watanabe2009}. \LCalc{} For plug-in or
gradient-trained predictors, $\kappa$ is the learner's own generalization
coefficient, which need not equal $\lambda$ in singular models. SLT
therefore describes Bayes codes directly and other codes only through a
separate comparison of $\kappa$ with $\lambda$.
\item[Priors.] \LCalc{} At attainable $n$, the prior-dependent terms of
\eqref{eq:gain} can dominate the logarithmic term, so every interpretation
in terms of $\lambda$ is a hypothesis to be tested.
\item[Tests.] \LCalc{} A Bayes factor is an e-value against the control's
prior predictive, and its sequential version is a test martingale under
that null. Universal inference gives a prior-free e-value against a
composite control family only if the family's maximized likelihood can be
bounded from above with certainty. That is possible for convex families,
such as logistic readouts on fixed features, where duality gaps certify
the maximum. For neural control families an optimizer supplies only a
lower bound, and no guarantee is claimed. Products of e-values are used only where conditional validity holds for the full history, as for sequential likelihood ratios; otherwise evidence is combined by prespecified weighted averages. Details: Appendix~\ref{app:evidence}.
\end{description}

\subsection{Stage 1 validation: the decision experiment}

The first six weeks test whether Stage~1 adds value, with manufactured
explanations as the adversary.
\begin{description}[style=nextline,leftmargin=1.5em]
\item[Exact examples, weeks 1--2.] Gaussian and small singular models in
which Bayes redundancy, batching effects and optimizer effects can be
separated exactly, reusing the conjugate machinery of version 5.
\item[Ground truth and manufactured alignments, weeks 3--4.] Compiled
transformers \cite{tracr2023}, semi-synthetic models with known circuits
\cite{interpbench2024} and the causal-variable tasks of \cite{mib2025}.
Manufactured alignments are built to pass existing tests: expressive maps
on untrained networks that reach high held-out accuracy
\cite{sutter2025}, alignments at wrong loci that pass independent and
identically distributed (i.i.d.)\ splits, and alignments of wrong
algorithms onto networks that compute correlated variables.
\item[Strong baselines on unseen tasks, weeks 5--6.] I.i.d.\ held-out IIA;
distribution-shifted held-out IIA, for example with disjoint vocabularies
\cite{sutter2025}; the causal abstraction error with interventions on
unmapped variables \cite{meloux2026}; online probing codes without controls
\cite{voita2020}; and accuracy against map size \cite{pimentel2020pareto}.
The metric diagnostics D1--D3 of Appendix~\ref{app:topo} are computed for
every alignment.
\end{description}

\begin{hypbox}{H-A1 (rejecting convincing false explanations).}
Stage~1 gains, computed against untrained-network and wrong-locus controls
on fixed targets, reject manufactured alignments that pass
distribution-shifted held-out tests and the causal abstraction error, while
retaining valid explanations. A valid explanation counts as retained if
any member of its equivalence class of abstractions is retained.
\end{hypbox}

\begin{hypbox}{H-A2 (training trajectories).}
Along training checkpoints of a network that learns the task, the gain
against the untrained-network control rises with task competence, while
i.i.d.\ held-out IIA with expressive maps stays near its maximum throughout,
as Sutter et al.\ observed \cite{sutter2025}.
\end{hypbox}

\begin{hypbox}{H-A3 (SLT on known ground).}
In map families whose learning coefficients are known, per-sample Bayes
codes---or block codes corrected by the batching factor---have logarithmic
gain coefficients that match $\Delta\lambda$ within stated error, and the
prior-dependent terms are small enough at attainable $n$ for the match to
be observable.
\end{hypbox}

\begin{hypbox}{H-A4 (a complete protocol on the benchmark).}
On the benchmark of \cite{meloux2026}, a protocol that combines Stage~1
gains with interventions on unmapped variables discriminates valid from
invalid abstractions better than either component alone.
\end{hypbox}

\subsection{Stage 2: explanations as codes for network behaviour}
\label{sec:a1a2}

\LTest{} Stage~2 codes the network's own responses, so that competing
explanations predict the same observations. An intervention protocol fixes
a distribution over inputs and interventions: interchange interventions at
the loci that explanations name, interventions on variables they leave
unmapped, and random-direction controls. The data are the network's output
distributions under these interventions or, following \cite{ma2026},
distributions over future operations. Each explanatory model receives a
stated set of information: the input, a description of the intervention,
and nothing about the network's internals beyond what the explanation
posits. The competitors are causal abstractions with their alignment
families, data-driven interfaces of the kind compared in \cite{ma2026},
and capacity-matched black-box predictors of intervened outputs. Models
are scored by the prequential code length of the network's responses;
e-values against tractable null families give guaranteed rejections; and
the complexity of explanation families enters through their learning
coefficients where these are known.

\paragraph{What counts as one observation.} Before any comparison, Stage~2
fixes the unit of observation: a sampled output under an intervention, an
exactly available probability vector, or a distribution estimated from a
stated number of samples. Cross-entropy against a probability vector is an
expected coding loss, not the likelihood of an observation, and the choice
changes the meaning of $n$, of Bayesian complexity penalties and of
e-values. The intervention protocol is frozen before scoring, so that
candidate explanations cannot change the experiment on which they are
judged.

\paragraph{Independent justification.} Stage~2 asks whether explanations
compress the network's own behaviour, a question Stage~1 cannot answer, so
it does not inherit its justification from Stage~1. It proceeds only after
a two-week pilot on one ground-truth model shows that shared-response codes
separate a valid abstraction from a manufactured one.

\begin{hypbox}{H-A5 (explanations beat black boxes on shared data).}
On ground-truth models, valid causal abstractions code the network's
interventional responses more briefly than capacity-matched black-box
predictors and than manufactured explanations, and the differences survive
correction for batching and prior effects.
\end{hypbox}

\paragraph{Relation to forked futures.} Ma et al.\ compare Shared, Local,
Mixture and Distributed interfaces by prequential causal description length
under fidelity and capacity constraints. They report held-out gains on
Qwen2.5-1.5B and Llama-3-8B and add transplantation, mediation and a blind
model-organism test \cite{ma2026}. Stage~2 differs in three ways. It tests
researcher-specified causal abstractions, including deliberately
manufactured ones, against data-driven interfaces and black boxes; it
comes with a theory of what the codes measure; and it offers guaranteed
tests where they exist. Their results are treated as preprint claims, to be
reproduced where possible, and the authors are natural collaborators.

\paragraph{Applications.} If Stage~2 works on ground truth, it is applied
first to refusal, whose behaviour and interventions are comparatively
concrete \cite{arditi2024}, and then to deception and sandbagging probes
\cite{goldowskydill2025}. Sparse-autoencoder and transcoder features
\cite{mib2025} and components singled out by susceptibility clusters
\cite{baker2026,gordon2026} enter as candidate variables.

\begin{hypbox}{H-A6 (safety directions).}
For refusal and deception directions in open language models, explanation
codes separate directions that causally mediate the behaviour from
directions that only predict it, and their ranking agrees with causal
ablations more often than probe accuracy does. A direction that explains few bits relative to its controls offers little incremental evidence of mediation under the chosen learners and controls. Limited data or an unsuitable learner can also produce low gain, so such a result calls for scrutiny of monitors built on the direction rather than a verdict on it.
\end{hypbox}

\begin{hypbox}{H-A7 (a data-side check).}
For valid abstractions, observables on the aligned subspace---for example,
the expected aligned coordinate over a fixed probe set, as a function of
the weights---are markedly more susceptible to data shifts that change the
distribution of $V$ than to shifts that do not. For manufactured alignments
this selectivity is absent.
\end{hypbox}

\subsection{Proof targets}

\begin{keybox}{Theory targets.}
\begin{description}[style=nextline,leftmargin=1.2em]
\item[T1, decomposition on fixed targets.] A rigorous version of
\eqref{eq:gain} with randomized-response targets, explicit conditions, and
the fluctuation and prior terms. Status: follows from Watanabe's theorem
applied to each code; to be written out.
\item[T2, label-only control.] Relabelling-invariant priors give zero gain
against the label-only code. Status: proved (Appendix~\ref{app:codes}).
\item[T3, batching factor.] For learners with excess loss
$\kappa/m+o(1/m)$ and geometric blocks, the cumulative excess is $\kappa\frac{r-1}{\ln r}\log n+o(\log n)$, with an $O(1)$ remainder when the learning curve's own remainder is $O(m^{-1-\delta})$ for some $\delta>0$. Status: derived, and checked exactly
for the Gaussian mean (Appendix~\ref{app:batch}); the general statement is
routine to write out.
\item[T4, learner coefficients.] Conditions under which a learner's
coefficient $\kappa$ equals $\lambda$, and singular families in which
plug-in predictors have $\kappa\neq\lambda$. Status: open.
\item[T5, tests.] Validity of Bayes-factor e-values against a control's
prior predictive, of universal-inference e-values for control families
with certified likelihood bounds, and of sequential products with fresh
data. Status: follows from standard e-value theory
\cite{grunwald2024safe,wasserman2020}; to be written for this setting.
\item[T6, untrained-network control.] Conditions under which a
structureless injective encoding forces a larger complexity term than a
network that computes $V$. Status: hypothesis.
\item[T7, metric certificates.] Finite realizability by diffeomorphisms, the
stretch budget and its persistence form, interleaving and realizability
bounds, and the prior mass of distortion. Status: proved or sketched in
Appendix~\ref{app:topo-protocol}; the link between distortion and learning
coefficients is open.
\end{description}
\end{keybox}

\subsection{Rivals, neighbours and positioning}

The rivals are i.i.d.\ and distribution-shifted held-out IIA
\cite{sutter2025}; the causal abstraction error with interventions on
unmapped variables \cite{meloux2026}; online probing codes \cite{voita2020}
and control tasks \cite{hewittliang2019}; accuracy against complexity
\cite{pimentel2020pareto}; and prequential causal description length over
forked futures \cite{ma2026}. Neighbours include a category-theoretic
framework grounding mechanistic explanation in minimum description length
\cite{gauderis2026}, description-length views of sparse autoencoders
\cite{ayonrinde2024}, proof length as a measure of understanding
\cite{gross2024}, the use of loss-landscape degeneracy for interpretability
\cite{bushnaq2024}, ARC's surprise accounting \cite{hilton2024} and
susceptibilities \cite{baker2026}.

Project~A's distinctive contributions are four: adversarial validation
against manufactured explanations and the strongest current baselines; a
theory of what prequential scores measure, with SLT describing the
Bayesian part; Stage~2's comparison of researcher-specified abstractions,
data-driven interfaces and black boxes on shared data; and safety
applications with guaranteed tests where these exist.

\subsection{Gates and outputs}
\label{sec:gatesA}

\begin{description}[style=nextline,leftmargin=1.5em]
\item[G-A0, end of week 2.] The exact examples separate Bayes redundancy,
batching and optimizer effects, and reproduce the batching factor.
\item[G-A1, end of week 6 (mid-November 2026): the decision gate.] Does the
Stage~1 score reject convincing false explanations that the strongest
baselines accept, while retaining valid explanations up to equivalence?
If it passes, the method paper is written for ICML 2027, whose deadline is
estimated for late January. If it fails, Stage~1 stops as a paper project and the finding is shared as a report or post. Further work on Stage~1 then receives a bounded feasibility budget of at most three weeks and continues only with a specific new hypothesis and preliminary evidence for it. Stage~2 proceeds only on its own justification (Section~\ref{sec:a1a2}).
\item[G-A2, month 5 (February--March 2027).] Stage~2 on ground truth: H-A5 holds in at least two model families, with batching and prior corrections applied. If it fails, Stage~2 stops as a paper project and its budget returns to the reserve.
\item[G-A3, month 7 (May 2027).] The refusal application succeeds with
controls. If so, Stage~2 and the refusal study are submitted to NeurIPS
2027.
\end{description}
\noindent Outputs: the Stage~1 paper and library if G-A1 passes; the
Stage~2 paper; a theory note, if the theory strand goes beyond the batching factor and standard e-value constructions; and an Alignment Forum post at
each stage.

\section{Project B (gated): recoverability audits with a geometric statistic}
\label{sec:projB}

\subsection{The question}

Auditors need to know whether a capability that a model does not display
could be recovered cheaply. Elicitation can improve benign models as well
as sandbagging ones, so improvement alone is an unreliable detector
\cite{taylor2025}. For an arbitrary model the defensible conclusion is
operational: \emph{capability $C$ is recoverable from weights $w_0$ under
a stated adaptation protocol and budget}. ``Suppressed'' and ``absent''
are used only for organisms whose training history is known. A capability
can be cheap to acquire because it is simple, and an existing capability
can be hard to recover with a given adapter family. Project~B asks whether
one independently measured geometric statistic predicts recoverability
better than excess description length and calibrated uplift.

\subsection{Theory}

\begin{keybox}{Excess description length, corrected (internal calculation).}
EDL is $L^A(n)-n\,L_{\rm test}(\theta^A_n)$: the prequential code length of
the fine-tuning labels minus $n$ times the final model's test loss
\cite{donoway2026}. Under a Bayesian data model the prequential part splits
exactly,
\begin{equation}
\E\,L^A(n)=\E\,F_n+\mathrm{AC}_A(n),\qquad
\mathrm{AC}_A(n)=\sum_{i=1}^n\E\,\KL\bigl(q_i\,\|\,p^A_i\bigr)\ge0,
\label{eq:edlsplit}
\end{equation}
with $q_i$ the Bayes predictive after $i-1$ labels. Hence $\E\,\EDL^A$ is
the Bayes part $\E F_n-n\,\E L_{\rm test}(q_{n+1})$, which equals
$\lambda\log n-(m-1)\log\log n+O(1)$, plus $\mathrm{AC}_A(n)$, minus $n$
times the excess of the learner's final test loss over the Bayes
predictive's. Details: Appendix~\ref{app:edl}.
\end{keybox}

That elicitation corresponds to a small learning coefficient of the
unlocking adaptation, and teaching to a large one, is a hypothesis, not a
consequence: prior terms and learner coefficients can carry the
difference. Optimizer effects on EDL lie in the algorithmic terms by
construction. The substantive question is whether the algorithmic excess,
estimated against an independent Bayes estimate, predicts
optimizer-dependent behaviour such as that reported in \cite{brown2026}.

\paragraph{Local free-energy depth.} \LCalc{} Let $U_{\rm on}$ and
$U_{\rm off}$ be the regions of adaptation space near $w_0$ in which the
capability is and is not displayed, defined by a fixed evaluation
threshold. Under the local tempered posterior on neutral data, the depth
$D=F_{\rm on}-F_{\rm off}$ is the log-odds against the capable region at a
fixed reference scale. Adding $k$ demonstrations while keeping the neutral
data multiplies the posterior by $e^{-\beta S_k}$, where $S_k$ is the summed
loss on the demonstrations. The log-odds of the capable region therefore
change exactly by
\begin{equation}
\log\E_{\rm on}\bigl[e^{-\beta S_k}\bigr]-\log\E_{\rm off}\bigl[e^{-\beta S_k}\bigr],
\label{eq:logodds}
\end{equation}
with expectations over the unmodified local posterior restricted to each
region (Appendix~\ref{app:flip}). Along the continuous population-loss tilt $p_t\propto p_0e^{-\beta tL_C}$, the initial slope in $t$ is $\beta\,(\E_{\rm off}[L_C]-\E_{\rm on}[L_C])$, which involves the capability loss alone; actual demonstrations are handled by the exact identity. The replacement formula of version 6 agrees only
when both regions fit the neutral data equally well. Computing
\eqref{eq:logodds} by reweighting, or by thermodynamic integration over $k$
when the weights degenerate, gives the flip point without extrapolating a
derivative.

\paragraph{Why depth is hard to measure.} Stable chains and reproducible
loss estimates do not establish relative basin masses: every chain can
remain in one region and agree with the others while missing the capable
region. Depth estimates therefore require defined regions, overlap along
the integration path and a test that would detect missed regions.
Project~B validates depth where masses are known before using it anywhere
else.

\subsection{Stage B1: validation where masses are exact}

Stage~B1 uses small models in which the posterior masses of the capable and
incapable regions can be computed by quadrature or by long, heavily
tempered reference runs. These include low-dimensional parametric models
with two basins and small networks trained with and without a password
lock. Stochastic-gradient estimates of $D$ are compared with the
references, including designed cases in which the capable region is hard
to reach.

\begin{hypbox}{H-B1 (depth is measurable).}
In small models with exact or reference masses, stochastic-gradient
estimates of $D$ agree with the references within a stated tolerance, and
the missed-region test flags every designed case in which chains fail to
reach the capable region.
\end{hypbox}

\subsection{Stage B2: matched organisms}

Stage~B2 uses matched organisms on open bases with one to two billion
parameters: password-locked models \cite{greenblatt2024}, imitation-trained
sandbaggers with distractor behaviours \cite{taylor2025}, benign models
fine-tuned identically without the lock, and never-trained controls built
by withholding synthetic facts. Four statistics are compared: calibrated
uplift (improvement on the target minus improvement on a matched control
task, benchmarked on benign models), EDL with the corrected definition,
noise injection \cite{tice2025} and depth $D$.

\begin{hypbox}{H-B2 (depth adds information).}
In matched organisms, $D$ improves prediction of recoverability---whether the stated adaptation procedure reaches the stated performance threshold within its budget---beyond what EDL and calibrated uplift predict together, at matched compute.
\end{hypbox}

\begin{hypbox}{H-B3 (the cost of unlocking).}
The flip point computed from \eqref{eq:logodds} with local samples of the
unmodified model agrees within a factor of two with the number of
demonstrations at which few-shot fine-tuning recovers the capability.
Large disagreements identify cases in which fine-tuning does not behave
like posterior updating.
\end{hypbox}

\paragraph{Two outcomes, kept apart.} Recoverability is the primary outcome of
the first Project~B paper. Whether an organism was built with a targeted
suppression or hidden-behaviour mechanism is a different target: a benign
model can contain an easily elicited capability, and predicting its
recovery correctly is not a false positive. Discrimination between training
histories is a separate, secondary analysis.

\paragraph{Audits with error guarantees.} For that secondary analysis, each
statistic can be converted into an e-value against the null that the
model's history involves no targeted suppression, by conformal calibration
against a library of models built without it \cite{vovk2005}. The
guarantee requires the audited model, under the null, to be exchangeable
with the library's models, which the organism design must make plausible.
E-values from different statistics or rounds are combined by a prespecified
weighted average; products are used only where conditional validity for
the full audit history has been shown (Appendix~\ref{app:evidence}).

\subsection{Stretch goals}
\label{sec:emcase}

If Stage~B2 succeeds, four extensions follow.
\begin{itemize}[leftmargin=*]
\item \textbf{Scale and realism.} Organisms with seven to eight billion
parameters, including backdoored ``sleeper agent'' models and a model with
a concealed objective, both re-created on open models following
\cite{hubinger2024} and \cite{marks2025}.
\item \textbf{Susceptibilities.} The covariances behind the initial slope
of \eqref{eq:logodds} are susceptibilities to added data
(Appendix~\ref{app:lr}); they also map the components that carry a
suppression \cite{baker2026}.
\item \textbf{Rare outputs.} Nearby models drawn from the local posterior
serve as proposals for importance sampling of the original model's rare
harmful outputs, with weights computed exactly from token log-probabilities
and each sample scored under every mixture component. H-B4: this reaches a
given relative error with fewer samples than isotropic weight noise and is
competitive with activation-steered proposals \cite{angell2026,wuhilton2025}.
\item \textbf{Emergent misalignment.} Broad misbehaviour after narrow
harmful fine-tuning \cite{betley2026} serves as a case study, using
existing small organisms \cite{turner2025}; general misalignment is
expected to be cheaper to recover than narrow misalignment
\cite{soligo2026,brown2026}.
\end{itemize}
Patterning suggests a defensive converse: choose unlearning data that
increases $D$, so that removed capabilities are harder to recover
\cite{wangmurfet2026}.

\subsection{Gates and outputs}

\begin{description}[style=nextline,leftmargin=1.5em]
\item[G-B1, month 6.] H-B1 holds. If it fails, depth is not used and the validation failure is shared as a report. Project~B continues only if a specific new hypothesis about recoverability passes a two-week feasibility check; otherwise its budget returns to the reserve.
\item[G-B2, month 9.] H-B2 holds on matched organisms. If so, the recoverability paper is submitted to ICLR 2028, with the deadline expected in late September or early October 2027. If not, the paper is not pursued and the stretch goals are cancelled.
\item[G-B3, month 10.] Stretch goals begin if G-B2 passes.
\end{description}
\noindent Stages~B1 and B2 need a few hundred A100-hours; the stretch goals
up to about 2,000.

\section{The dynamics note: learning coefficients predict selection, not passage}
\label{sec:p0}

This short note (P0) closes the question of the earlier versions in
public, engages a published claim with a checkable correction, and
states the lesson that motivates version 6. It reuses completed work;
the budget is 0.75 person-months, and the Alignment Forum version comes first (weeks 6--8), with the arXiv version after the ICML deadline (month 4).

\begin{enumerate}[leftmargin=*]
\item \textbf{Scoped impossibility.} \LCalc{} No function of the saddle
invariants $(\lambda_s,m_s)$, or of any population-loss statistic,
determines saddle-passage statistics once teacher, gradient covariance,
initialization scales or metric may vary. The four controls from
version 5 establish this (Appendix~\ref{app:p0}). The statement does not
cover teacher-dependent finite-temperature quantities under a fixed
optimizer, and $\lambda$ does enter activated escape rates through the
masses of wells.
\item \textbf{Flat sets.} \LCalc{} and \LNum{} The published passage law
$T\propto\eps^\lambda$ for diffusion through a flat region
\cite{hennick2026} fails even for isotropic diffusion in a reflecting
tube. The passage time scales as $\eps^{\lambda-\chi_{\rm exit}}$, up to
logarithms, where $\chi_{\rm exit}$ is the scaling exponent of the exit
cross-section. Numerical checks confirm constant, logarithmic and
$\eps^{-1/4}$ behaviour in three cases where $\eps^\lambda$ predicts
decay (Appendix~\ref{app:p0}).
\item \textbf{Discrete chain.} \LNum{} The continuous approximation is
accurate for small steps. The strong-noise power law never appears in
discrete SGD. With label noise about 0.9\% of paths cross zero and
never exit. Exit time is not a function of $\eta/b$ alone.
\item \textbf{The positive lesson.} $\lambda$ governs posterior masses and
code lengths, that is, selection and evidence, not passage times. A
basin-selection account of grokking makes the same distinction in a
different setting \cite{cullen2026}.
\end{enumerate}
\noindent A courtesy copy goes to the authors of \cite{hennick2026}
before posting.

\section{Tooling}
\label{sec:tools}

The \texttt{devinterp} library is actively maintained. The first step is
an issue, not a pull request, proposing (i) regression tests against
exact learning coefficients (monomial germs and reduced-rank regression
\cite{aoyagiwatanabe2005}) and (ii) a code-length utility with per-sample and block prequential estimators, the batching-factor correction, estimators of the widely applicable Bayesian information criterion \cite{watanabe2013} and the thermometer. A pull request
follows in months 6--9 if the maintainers want it. Project~A's evaluation library is released with its preprint. A third contribution is a calibration suite for susceptibility estimates: exact posterior covariances in conjugate Gaussian models and, by quadrature, in two-dimensional singular potentials, against which stochastic-gradient estimates are checked across step sizes and chain lengths. It could become a short joint paper with Resolution. Budget: 0.5 person-months before month 12 and 0.25 afterwards.

\section{Schedule and effort}
\label{sec:schedule}

The plan assumes a full-time commitment: about twelve person-months before
month 12 and eighteen in total. At most two projects are in active
execution at once, besides small tooling and writing tasks, and the
reserve is released by gate outcomes rather than by drift.

\begin{center}
\begin{tabular}{@{}P{2.4cm}P{8.1cm}P{4.4cm}@{}}
\toprule
Period & Primary work & Public milestone \\
\midrule
Weeks 1--2 (Oct 2026) & Exact examples and G-A0; draft of the dynamics
note; \texttt{devinterp} issue. & --- \\
Weeks 3--6 (Oct--Nov) & Ground truth, manufactured alignments and strong
baselines; decision gate G-A1. & Dynamics note on the Alignment Forum
(weeks 6--8). \\
Weeks 7--16 (Nov 2026--Jan 2027) & If G-A1 passes: Stage~1 paper and library. Otherwise: report, bounded feasibility budget, and the Stage~2 pilot. Stage~B1 set-up. & Stage~1 paper (ICML 2027) if G-A1 passes. \\
Months 4--7 (Feb--May) & Stage~2 on ground truth (G-A2); refusal study
(G-A3); Stage~B1 validation (G-B1, month 6). & Dynamics note on arXiv
(month 4); NeurIPS 2027 submission. \\
Months 8--12 (Jun--Oct) & Stage~B2 matched organisms (G-B2, month 9);
deception probes; stretch goals from month 10. & ICLR 2028 submissions;
NeurIPS decisions; posts. \\
Months 12--18 & Revisions, stretch goals, journal versions, talks and
applications. & Record complete by month 12. \\
\bottomrule
\end{tabular}
\end{center}

\begin{center}
\begin{tabular}{@{}lccc@{}}
\toprule
Item (person-months) & By month 12 & Months 12--18 & Total \\
\midrule
Dynamics note & 0.75 & 0.25 & 1.0 \\
Project A, Stage 1 and theory strand & 3.0 & 0.5 & 3.5 \\
Project A, Stage 2 and applications & 3.5 & 1.0 & 4.5 \\
Project B & 3.25 & 2.0 & 5.25 \\
Tooling & 0.5 & 0.25 & 0.75 \\
Outreach, collaboration and applications & 0.5 & 1.0 & 1.5 \\
Reserve & 0.5 & 1.0 & 1.5 \\
\midrule
Total & 12.0 & 6.0 & 18.0 \\
\bottomrule
\end{tabular}
\end{center}

\paragraph{Stop rules and the reserve.} A line of work that fails its gate stops as a paper project, and its remaining budget returns to the reserve. The reserve funds only lines that have passed their gates, or new hypotheses with preliminary evidence from a bounded feasibility budget: at most three weeks after G-A1 and two weeks after G-B1 or G-B2. Stage~2 must pass its own pilot whatever the outcome of G-A1.

\paragraph{Readiness indicators.} Instead of publication probabilities,
the plan tracks what each output depends on and how close it is.

\begin{center}
\begin{tabular}{@{}P{2.9cm}P{2.4cm}P{9.6cm}@{}}
\toprule
Output & Depends on & Readiness indicators \\
\midrule
Dynamics note & --- & Draft complete; flat-set correction checked against
the original assumptions; courtesy copy sent. \\ \addlinespace
Theory note & G-A0, T4 & Results beyond the batching factor and standard e-value constructions, such as learner coefficients in singular families. \\ \addlinespace
Stage~1 paper & G-A0, G-A1 & Share of manufactured alignments rejected
that the strongest baseline accepts; false-rejection rate for valid
explanations; effect sizes with intervals; tested library. \\ \addlinespace
Stage~2 paper & G-A2, G-A3 & Shared response datasets; capacity-matched
black boxes; corrected gains; refusal study with controls. \\ \addlinespace
Project~B paper & G-B1, G-B2 & Validation on exact masses; missed-region
test; matched organisms; effect sizes for H-B2. \\
\bottomrule
\end{tabular}
\end{center}

\paragraph{Contacts.} Month 2: Resolution, with the dynamics note, the
Project~A plan and a proposal for joint work on susceptibilities and
calibration. After G-A1: the authors of
\cite{ma2026,sutter2025,meloux2026,gauderis2026}. During the theory strand:
ARC and Gr\"unwald's group on e-values. After G-B1: the authors of the EDL
and noise-injection work.

\section{Outlook}
\label{sec:outlook}

Failed experiments are worth sharing, but not every outcome supports a
paper: establishing that familiar baselines work on the examples tried is
not automatically a significant negative result. The table lists what
would be shared under each outcome and how the rest of the plan responds.
Subjective publication probabilities are kept outside the plan and revised
at each gate.

\begin{center}
\begin{tabular}{@{}P{3.4cm}P{6.0cm}P{5.5cm}@{}}
\toprule
Outcome & Shared as & Effect on the rest \\
\midrule
G-A1 passes & Stage~1 paper (ICML 2027 or NeurIPS 2027) and library &
Stage~2 proceeds if its own pilot succeeds. \\ \addlinespace
G-A1 fails & Report or post; a paper only if the failure reveals something
unexpected about strong baselines or about the codes & Stage~1 stops;
bounded feasibility budget; Stage~2 needs its own justification. \\
\addlinespace
Theory strand & A note only if it goes beyond the batching factor and
standard e-value constructions & Otherwise the results serve as
foundations for the papers. \\ \addlinespace
G-A2 and G-A3 pass & Stage~2 paper with the refusal study (NeurIPS 2027) &
Deception probes follow. \\ \addlinespace
G-A2 fails & Report & Stage~2 stops; its budget returns to the reserve. \\
\addlinespace
G-B1 fails & Report of the validation failure & Project~B stops unless a
new hypothesis passes a two-week feasibility check. \\ \addlinespace
G-B2 passes & Recoverability paper (ICLR 2028) & Stretch goals. \\
\addlinespace
G-B2 fails & Report & Stretch goals are cancelled. \\
\bottomrule
\end{tabular}
\end{center}

\noindent Under unfavourable outcomes the record by month 12 may consist of
the dynamics note, reports and tooling, which is thin for the objective.
That is the risk the plan accepts in exchange for ambitious targets; the
stop rules limit the time it can cost. Under favourable outcomes the record
contains a method paper, a language-model-scale interpretability paper
with a safety application and a recoverability paper under review.

\section{Risks and responses}
\label{sec:risks}

\begin{description}[style=nextline,leftmargin=1.5em]
\item[R1: overlap or scoop.] Close work exists \cite{ma2026} and active
groups work on causal-abstraction metrics and description-length
frameworks \cite{meloux2026,gauderis2026}. Response: speed on the decision
experiment, differentiation through adversarial validation and theory, and
early contact.
\item[R2: the baselines already solve the problem.] Distribution-shifted
held-out tests or complete intervention protocols may expose every
manufactured alignment tried. Accepted: Stage~1 then stops as a paper project under the stop rule, and the finding is shared.
\item[R3: estimator artefacts.] Batching, learner and prior effects can
masquerade as geometry. Response: the theory strand, exact examples and
per-sample Bayes codes for every SLT comparison.
\item[R4: the SLT hypotheses fail.] H-A3 or the link between elicitation
and small learning coefficients may not hold. Accepted and reported; the
main statistics do not depend on them.
\item[R5: no guaranteed tests for neural controls.] Certified likelihood
bounds exist only for tractable families. Response: guarantees are claimed
only there; other comparisons are reported as relative fit.
\item[R6: depth is not measurable.] Response: G-B1 on exact masses before
any use at scale.
\item[R7: unrealistic organisms.] Response: imitation-trained sandbaggers,
matched benign controls, and treating success as necessary rather than
sufficient evidence for real audits.
\item[R8: overcommitment.] Response: at most two projects in active
execution, and gate-driven release of the reserve.
\item[R9: engineering at scale by one person.] Response: stage every
experiment at one to two billion parameters, reuse open organisms and
libraries, and seek a collaborator for Project~B's infrastructure.
\item[R10: the ICML date moves.] The deadline is an estimate. An earlier
date means a Stage~1 paper built on the ground-truth results alone; a
later one adds the benchmark study.
\item[R11: dual use.] The tools are defensive. No misaligned adapters are
released beyond those already public.
\end{description}

\appendix

\section{Technical material for the dynamics note}
\label{app:p0}
\renewcommand{\theequation}{A.\arabic{equation}}
\setcounter{equation}{0}

\subsection{Four controls from version 5}

\LCalc{} Each row holds the model-intrinsic invariants fixed and changes
one ingredient of the training process. The derivations are in version 5
(Appendices A and B and Section 6.2).

\begin{center}
\begin{tabular}{@{}P{2.3cm}P{3.8cm}P{3.6cm}P{4.8cm}@{}}
\toprule
Control & Held fixed & Changed & Consequence for exit \\
\midrule
Teacher & Self-divergence and $(\lambda_s,m_s)=(1,L)$ & Realizable
teacher, cancelling versus non-cancelling & Exit time
$\propto u_0^{-(2L-2)}$ versus $u_0^{-(L-2)}$ \\ \addlinespace
Gradient covariance & Population loss, Fisher, self-divergence &
Gaussian versus Rademacher inputs & Positive versus exactly zero
minibatch noise \\ \addlinespace
Scale hierarchy & Self-divergence, teacher & Initialization
$(\eps,\eps,\sqrt\eps,1)$ versus equal small layers &
$\tfrac12\eps^{-1/2}\log(1/\eps)$ versus $\eps^{-(r-2)}$ \\ \addlinespace
Metric & Population loss, self-divergence & Euclidean versus natural
gradient & Polynomial delay versus bounded time; fixed damping restores
the delay \\
\bottomrule
\end{tabular}
\end{center}

\subsection{Passage through flat sets}

\LKnown{} Hennick and De Baerdemacker argue that the time for isotropic
diffusion to pass through a small flat region scales as $\eps^\lambda$
divided by the capacitance of the escape set, which is treated as
independent of $\eps$ \cite{hennick2026}.

\LCalc{} Consider diffusion with coefficient $D$ in the tube
$W_\eps=\{K<\eps\}$ inside a fixed domain, reflecting on $\{K=\eps\}$ and
absorbed at an escape set. For thin tubes the Fick--Jacobs reduction
\cite{zwanzig1992,reguera2001} gives a one-dimensional problem along the
zero set, with cross-section $A(s)$ and enclosed volume
$V(s)=\int_0^sA(u)\,\dd u$. Starting at the centre, with reflection there
by symmetry,
\begin{equation}
T\approx\frac1D\int_0^{s_{\rm exit}}\frac{V(s)}{A(s)}\,\dd s .
\label{eq:fj}
\end{equation}
Three germs illustrate the result, with $D=1$ for the numbers.
\begin{enumerate}[leftmargin=*]
\item $K=y^2$ in the unit disk ($\lambda=\tfrac12$). Here $A=2\sqrt\eps$
and $V=2\sqrt\eps\,x$, so $T=\int_0^1x\,\dd x=1/(2D)$, independent of
$\eps$.
\item $K=x^2y^2$ in the unit disk ($\lambda=\tfrac12$, $m=2$). Here
$A(x)=2\min(x,\sqrt\eps/x)$, $V(x)=\sqrt\eps+2\sqrt\eps\log(x\eps^{-1/4})$
for $x>\eps^{1/4}$, and $T=\log(1/\eps)/(8D)+O(\sqrt\eps)$. The
multiplicity survives as a logarithm.
\item $K=x^4y^2$ in $[-1,1]^2$, absorbed only at $x=\pm1$ ($\lambda=\tfrac14$).
Here $A(x)=\min(2,2\sqrt\eps/x^2)$, $V(x)=4\eps^{1/4}-2\sqrt\eps/x$ for
$x>\eps^{1/4}$, and $T=\tfrac23\eps^{-1/4}/D-1/(2D)+O(\sqrt\eps)$: the
time grows as $\eps\to0$.
\end{enumerate}
In general the dominant volume scales as $\eps^\lambda$, up to
logarithms, while the exit cross-section scales as $\eps^{\chi_{\rm exit}}$.
The integral in \eqref{eq:fj} is dominated near the exit, giving
$T\propto\eps^{\lambda-\chi_{\rm exit}}$ up to logarithms. All three
cases have $\chi_{\rm exit}=\tfrac12$. The capacitance of an escape window
that is itself a cross-section of the tube scales with $\eps$; it is not
constant.

\LNum{} Finite-difference solutions of the reflected-diffusion
mean-first-passage problem (\texttt{tube\_mfpt.py}) agree with
\eqref{eq:fj}:

\begin{center}\small
\begin{tabular}{@{}ccccccc@{}}
\toprule
$\eps$ & (1) numerical & (2) numerical & (2) reduction & (3) numerical & (3) reduction & $\eps^\lambda$ law \\
\midrule
$10^{-2}$ & 0.500 & 0.620 & 0.576 & 1.835 & 1.608 & decreasing \\
$10^{-3}$ & 0.504 & 0.883 & 0.863 & 3.434 & 3.249 & decreasing \\
$10^{-4}$ & 0.500 & 1.150 & 1.151 & 6.253 & 6.167 & decreasing \\
$10^{-5}$ & --- & 1.430 & 1.439 & 11.359 & 11.355 & decreasing \\
\bottomrule
\end{tabular}
\end{center}

\noindent The value 0.504 reflects grid alignment. In case (3) the
fitted log--log slope is $-0.264$ overall and $-0.259$ on the last
decade, against the predicted $-1/4$; in case (2) the slope of $T$
against $\log(1/\eps)$ is $0.117$, against $1/8$. The calculation uses a
reflecting-tube caricature; whether it transfers to the exact
formulation of \cite{hennick2026} will be checked with the authors. It
is consistent with order-one motion along zero sets \cite{larsson2026}.

\subsection{The discrete balanced chain}

\LNum{} The balanced update of version 5 is
$u_{k+1}=u_k+\eta u_k^{L-1}[(s-u_k^L)A_k+C_k]$, with
$A_k=b^{-1}\sum x_i^2$ and $C_k=b^{-1}\sum q_ix_i$. We used $L=3$,
$s=1$, $w_0=10^{-3}$, $w_1=1/2$, Gaussian inputs and the clock $T=K\eta$
(\texttt{discrete\_sgd\_check.py}; 4000 paths per setting).

\begin{center}\small
\begin{tabular}{@{}lcccccc@{}}
\toprule
$\eta$ (with $b=1$, $\sigma_q=0$) & 0.01 & 0.1 & 1 & 3 & 10 & 100 \\
\midrule
This check & 9.19 & 9.75 & 13.98 & 20.70 & 38.3 & 179.0 \\
Standard error & 0.007 & 0.021 & 0.063 & 0.12 & 0.24 & 1.3 \\
Direction review & 9.19 & 9.79 & 14.05 & 20.7 & 38.7 & 182 \\
Diffusion, exact quadrature & 9.15 & 9.51 & 14.25 & --- & 186.9 & 14278 \\
\bottomrule
\end{tabular}
\end{center}

\noindent Without label noise $A_k\ge0$, so no path ever moves backward.
The diffusion is accurate up to $r=\eta/b\approx1$, and its strong-noise
power law never appears in discrete SGD.

With label noise ($\sigma_q=1$, $\eta=1$, $b=1$; \texttt{trapped\_check.py};
20{,}000 paths, $10^6$ steps), 179 paths did not exit ($0.90\%\pm0.07\%$)
and 212 crossed zero. Every non-exiting path sits near
$u\approx-1.0\times10^{-6}\approx-1/(\eta sk)$. For $u<0$ the drift
$\eta su^2$ pushes $u$ towards zero from below ever more slowly, and
returning above zero would require a noise kick of order $1/(\eta|u|)$.
The median exit time is 15 and the mean conditional on exit about 16.3;
the unconditional mean is infinite. The review's ``mean of order
$10^3$--$10^4$'' is inconsistent with its own censoring count: 13 of 2000
paths unfinished at $10^7$ steps imply a restricted mean of at least
$6.5\times10^4$. With $r=\eta/b=0.1$ fixed, $(\eta,b)=(0.1,1)$, $(1,10)$
and $(10,100)$ give $T=9.77$, $11.61$ and $29.71$ (review: 11.8 and 29.6),
so exit time is not a function of $r$ alone.

\section{Code-length algebra}
\label{app:codes}
\renewcommand{\theequation}{B.\arabic{equation}}
\setcounter{equation}{0}

\subsection{Prequential and mixture codes coincide for Bayes predictors}

\LKnown{} With the posterior $\pi(\cdot\mid D_{i-1})$ and predictive
$q_i(y\mid x)=\int p(y\mid x,w)\,\pi(w\mid D_{i-1})\,\dd w$, the chain rule
gives $q(y^n\mid x^n)=\prod_iq_i(y_i\mid x_i)$, so
$F_n=\sum_i-\log q_i(y_i\mid x_i)$ exactly.

\subsection{Asymptotics}

\LKnown{} Expansion \eqref{eq:watanabe} holds for analytic models on
compact parameter sets, with priors positive and smooth near the optimal
set, under realizability or a relatively finite variance condition
\cite{watanabe2009,watanabe2018}. For regular models $\lambda=d/2$ and
$m=1$. The Bayes training loss is $L_n(w_0)+O_p(1/n)$, so the residual
code length under the final predictive is $nL_n(w_0)+O_p(1)$.

\subsection{The gain on fixed targets}

\LCalc{} Apply \eqref{eq:watanabe} to the Bayes codes of a configuration
and a control on the same ordered targets, with optima $\theta$ and
$\theta_C$:
\begin{equation}
G(n)=n\bigl[L^C_n(\theta_C)-L_n(\theta)\bigr]+\Delta\lambda\log n
-\Delta(m-1)\log\log n+O_p(1).
\end{equation}
Write $n[L^C_n(\theta_C)-L_n(\theta)]=n\Delta\ell+R_n$, where $R_n$ is the
difference of centred empirical losses; in general $R_n=O_p(\sqrt n)$. With
randomized response, each target is replaced with known probability
$\eps_{\rm e}$ by a uniformly random label and the ordinary softmax
likelihood is used. If a configuration can represent the resulting
conditional law, it does so at finite logits, so the model is realizable
and analytic. Two configurations that both fit exactly make identical
predictions on the same noisy labels, so $\Delta\ell=0$ and $R_n=0$.
Whether a frozen network with a given map family can attain the target
probabilities is an empirical question, checked by calibration. The
$O_p(1)$ terms include the log prior density near the optimal set; they do
not vanish and differ between configurations with equal $\lambda$.

\subsection{The label-only control gains nothing from relabelling-invariant priors}

\LCalc{} Fix the evaluation pairs. A map family and prior induce a
distribution $Q$ on labellings $y^n\in\mathcal Y^n$. Suppose $Q$ is
invariant under every permutation of $\mathcal Y$ applied independently at
each pair. These permutations act transitively on $\mathcal Y^n$, so $Q$ is
uniform and $-\log Q(y^n)=n\log|\mathcal Y|$ for every labelling: the Bayes
code equals the uniform code, and the gain against the label-only code,
which has no access to the algorithm, is zero. The minimax-regret statement
of Section~\ref{sec:primer} is the worst-case analogue
\cite{shtarkov1987}. Practical map families are not invariant and carry
inductive biases of their own, which is why the untrained-network control
is needed: it measures what the learner achieves without the network's
help. Measuring the information a representation makes usable to a
restricted family is familiar from probing
\cite{pimentel2020info,xu2020,voita2020}.

\subsection{The batching factor}
\label{app:batch}

\LCalc{} Suppose a learner's expected excess log loss after $m$ examples is
$\kappa/m+o(1/m)$. Code the data in blocks with boundaries
$n_j=\lceil r^j\rceil$, using for each block the predictor trained on all
data before it. Block $j$ then contributes
$(n_{j+1}-n_j)\,\kappa/n_j+o(1)=\kappa(r-1)+o(1)$, and there are about
$\ln n/\ln r$ blocks up to $n$, so
\begin{equation}
\text{cumulative excess}=\kappa\,\frac{r-1}{\ln r}\,\ln n+o(\ln n).
\end{equation}
The remainder is $O(1)$ if the learning curve's remainder is $O(m^{-1-\delta})$ for some $\delta>0$; a term of order $1/(m\log m)$, such as the multiplicity term of a singular Bayes learning curve, accumulates into $\log\log n$. Per-sample updating ($r\to1$) gives $\kappa\ln n$. For the Bayes predictive,
$\kappa=\lambda$ \cite{watanabe2009}, so a Bayes predictor frozen within
doubling blocks has coefficient $\lambda/\ln2$. For plug-in predictors in
regular models, $\kappa=d/2$; for plug-in or gradient-trained predictors in
singular models, $\kappa$ need not equal $\lambda$.

\LNum{} For a Gaussian mean with known unit variance, the expected excess
of each predictor is available in closed form: $1/(2m)$ for the sample
mean, and an exact Kullback--Leibler expression for the Bayes predictive
with prior $N(0,1)$ and true mean $0.3$. Summing exactly over block
schemes and fitting against $\ln n$ at block boundaries between $2^8$ and
$2^{22}$ gives (\texttt{batching\_check.py}):

\begin{center}
\begin{tabular}{@{}lccc@{}}
\toprule
Scheme & Growth factor $r$ & Fitted coefficient & $\tfrac12\,(r-1)/\ln r$ \\
\midrule
Plug-in, per sample & 1 & 0.5001 & 0.5000 \\
Plug-in & 1.5 & 0.6166 & 0.6166 \\
Plug-in, doubling & 2 & 0.7213 & 0.7213 \\
Plug-in & 4 & 1.0820 & 1.0820 \\
Bayes, per sample & 1 & 0.4998 & 0.5000 \\
Bayes, doubling & 2 & 0.7208 & 0.7213 \\
\bottomrule
\end{tabular}
\end{center}

\noindent Logarithmic slopes of block codes must therefore be divided by
$(r-1)/\ln r$ before comparison with $\lambda$, and even then they estimate
the learner's coefficient $\kappa$ rather than $\lambda$.

\subsection{Excess description length}
\label{app:edl}

\LCalc{} EDL is defined as $\EDL^A(n)=L^A(n)-n\,L_{\rm test}(\theta^A_n)$,
with $L_{\rm test}(\theta)=\E_{(x,y)\sim\mathcal D}[-\log p_\theta(y\mid x)]$
\cite{donoway2026}. Let $\theta$ be drawn from the prior, inputs be
independent, and $y_i\mid x_i,\theta\sim p(\cdot\mid x_i,\theta)$. Then, given
the first $i-1$ labels and $x_i$, the label $y_i$ has law $q_i$, the Bayes
predictive, and for any predictor $p^A_i$ built from the same information
\begin{equation}
\E\bigl[-\log p^A_i(y_i\mid x_i)\bigr]
=\E\bigl[-\log q_i(y_i\mid x_i)\bigr]+\E\,\KL\bigl(q_i\,\|\,p^A_i\bigr).
\end{equation}
Summing gives \eqref{eq:edlsplit}, and therefore
\[
\E\,\EDL^A=\bigl[\E F_n-n\,\E L_{\rm test}(q_{n+1})\bigr]
+\mathrm{AC}_A(n)-n\bigl[\E L_{\rm test}(\theta^A_n)-\E L_{\rm test}(q_{n+1})\bigr].
\]
For a fixed teacher under the conditions of \eqref{eq:watanabe},
$\E F_n=nL(w_0)+\lambda\log n-(m-1)\log\log n+O(1)$ and
$\E L_{\rm test}(q_{n+1})=L(w_0)+\lambda/n+o(1/n)$, so the Bayes part is
$\lambda\log n-(m-1)\log\log n+O(1)$. The algorithmic term is nonnegative
under the Bayesian data model; the final-loss correction can have either
sign. The expected nonnegativity of EDL itself requires conditions on the
learner \cite{donoway2026}.

\subsection{How demonstrations move the local posterior}
\label{app:flip}

\LCalc{} Let $p_0$ be the local tempered posterior on neutral data. Adding
demonstrations $z_1,\dots,z_k$ with summed loss
$S_k(w)=\sum_{j\le k}\ell_C(w;z_j)$ gives $p_k\propto p_0\,e^{-\beta S_k}$.
Hence, exactly,
\begin{equation}
\frac{P_k(U_{\rm on})}{P_k(U_{\rm off})}
=\frac{P_0(U_{\rm on})\,\E_{\rm on}[e^{-\beta S_k}]}
{P_0(U_{\rm off})\,\E_{\rm off}[e^{-\beta S_k}]},
\end{equation}
with expectations under $p_0$ restricted to each region. This is
\eqref{eq:logodds}, a free-energy perturbation identity. An integer number of demonstrations does not define a derivative. Along the continuous population-loss tilt $p_t\propto p_0\,e^{-\beta tL_C}$, with $L_C$ the expected loss per demonstration, the log-odds have derivative $\beta\,(\E_{\rm off}[L_C]-\E_{\rm on}[L_C])$ at $t=0$; the finite identity above is what applies to actual demonstrations. Replacing neutral data by capability data at fixed
total size instead changes the log-odds at rate
$n\beta\,(\E_{\rm off}[\Delta L]-\E_{\rm on}[\Delta L])$ per unit mixture
weight, with $\Delta L=L_C-L_{\rm neutral}$. The two agree only if
$\E_{\rm on}[L_{\rm neutral}]=\E_{\rm off}[L_{\rm neutral}]$. Reweighting
estimates of \eqref{eq:logodds} degrade when the weights $e^{-\beta S_k}$
concentrate on few samples; thermodynamic integration over a sequence of
intermediate $k$ then replaces the single step.

\subsection{E-values from codes}
\label{app:evidence}

\LKnown{} and \LCalc{} (i) Let $q_H$ and $q_C$ be mixture codes. Since
$\sum_{y^n}q_H(y^n)=1$, we have $\E_{q_C}[q_H/q_C]=1$: the Bayes factor is an
e-value against the simple null $q_C$. Evaluated sequentially it is a
nonnegative martingale under that null, so by Ville's inequality the
probability that it ever exceeds $1/\alpha$ is at most $\alpha$.
(ii) For a composite null $\{p_{C,\theta}\}$, universal inference uses
$E=q^{\rm preq}(y^n)/\sup_\theta p_{C,\theta}(y^n)$, which satisfies
$\E_{p_{C,\theta}}[E]\le1$ for every $\theta$ \cite{wasserman2020}. The
supremum must be computed or bounded from above. Replacing it by an
optimizer's best value, which is a lower bound, inflates $E$ and voids the
guarantee. For convex control families, such as logistic or Gaussian
readouts on fixed features, the dual objective certifies an upper bound.
The logarithm of $E$ equals the gain minus the regret of the control's
code relative to its maximum-likelihood fit. (iii) If the control family
contains the law of the targets, every valid e-value against it has
expectation at most one under that law, so none grows except with small
probability. (iv) Products across rounds require $\E[E_t\mid\mathcal F_{t-1}]\le1$ under the null \cite{grunwald2024safe}. Sequential likelihood ratios satisfy this by construction. Repeated audits of the same model need not, even with fresh data and fresh reference models: a benign model with idiosyncratic features can score high in every round, so marginal calibration can hold in each round while conditional calibration fails. Without a proof of conditional validity, evidence is combined by a prespecified weighted average of e-values, which is always valid, or by a separately validated sequential procedure.

\subsection{Linear response of the local posterior}
\label{app:lr}

\LCalc{} Susceptibilities as defined in \cite{baker2026} describe a shift of
the data distribution at fixed size. With $L_h=(1-h)L_q+hL_{q'}$,
differentiating $\E_h[\phi]=Z_h^{-1}\int\phi\,e^{-n\beta L_h}\varphi$ gives
$\dd\E_h[\phi]/\dd h=-n\beta\,\mathrm{Cov}_h(\phi,L_{q'}-L_q)$, which is
\eqref{eq:suscept}. Tilting by the capability population loss at fixed neutral size, $p_t\propto p_0e^{-\beta tL_C}$, instead gives $\dd\E_t[\phi]/\dd t=-\beta\,\mathrm{Cov}_t(\phi,L_C)$.
The two differ by the covariance with the neutral loss. The initial slope
of \eqref{eq:logodds} is the second form applied to the indicator of the
capable region; both forms can be estimated from the same local samples.

\section{Calibration tools carried over}
\label{app:calib}
\renewcommand{\theequation}{C.\arabic{equation}}
\setcounter{equation}{0}

\subsection{A loss-fluctuation thermometer}

\LCalc{} Under a Gibbs law proportional to $e^{-K/T}$ near a minimum of
multiplicity one, the excess loss is asymptotically Gamma-distributed
with shape $\lambda$ and scale $T$. Hence $\hat T=\Var/\text{mean}$ and
$\hat\lambda=\text{mean}^2/\Var$.

\LNum{} Exact two-dimensional quadrature (\texttt{thermometer\_check.py})
recovers $T$ and $\lambda$ to machine precision for $x^2+y^4$
($\lambda=3/4$) and $x^2+y^2$ ($\lambda=1$) at $T=10^{-2},10^{-3},10^{-4}$.
For $x^2y^2$ ($\lambda=1/2$, $m=2$) it gives $\hat T/T=0.93,0.97,0.98$ and
$\hat\lambda=0.37,0.40,0.42$: the logarithmic bias expected from the
multiplicity.

\paragraph{Caveat.} \LCalc{} For SGD whose gradient noise is shaped like
the curvature, the stationary covariance near a minimum solves
$H\Sigma+\Sigma H=(\eta/B)H$ \cite{mandt2017}, so $\Sigma=(\eta/2B)I$.
That law is not Gibbs, and the thermometer then returns half the
participation ratio $(\operatorname{tr}H)^2/(2\operatorname{tr}H^2)$ instead
of $\lambda$. The plan therefore applies the thermometer only to samplers
whose target is Gibbs by construction, as a check on step-size bias.

\subsection{Identifiability of logarithmic terms}

\LNum{} After removing a constant, the regressors $(\log n)/n$ and $1/n$
have correlation $0.999$ over one decade of $n$: the variance inflation
factor is 568 and the standard error of the $\log n$ coefficient is
inflated 23.8-fold. Over two and three decades the factors are 348
(18.7-fold) and 401 (20.0-fold). For $T\log(1/T)$ against $T$ they are
568 (23.8-fold), 153 (12.4-fold) and 44 (6.6-fold)
(\texttt{identifiability.py}). Consequences: the $(\log n)/n$ signature
proposed for the emergent-misalignment boundary cannot be separated from
an $O(1)$ free-energy offset; multiplicity terms in code-length curves are
not claimed; and Project~B estimates flip points and occupancies rather
than fitting logarithmic coefficients.

\section{Topological and metric certificates for alignment maps}
\label{app:topo}
\renewcommand{\theequation}{\thesection.\arabic{equation}}
\setcounter{equation}{0}
\providecommand{\Lip}{\operatorname{Lip}}
\providecommand{\Dist}{\operatorname{Dist}}
\providecommand{\Dgm}{\operatorname{Dgm}}
\providecommand{\VR}{\mathrm{VR}}
\providecommand{\sep}{\operatorname{sep}}

This appendix asks whether persistent homology can test if an alignment map
supplies the structure that an interchange-intervention analysis reports.
\LAssess{} A search on 12 October 2026 found no work that applies persistent
homology to causal-abstraction alignment maps; the results below are derived
here. Neighbouring work compares representations through their multiscale
topology \cite{barannikov2022}, shows that trained networks simplify the
topology of class manifolds layer by layer \cite{naitzat2020}, quantifies how
the colour classes of a point cloud mingle \cite{cultrera2025chromatic}, and
bounds the persistent cost of Lipschitz maps \cite{gozzi2025}. A 2026 study of
activation-space topology during alignment fine-tuning \cite{phalign2026} uses
``alignment'' in the safety sense and is unrelated.

\paragraph{Summary.} Topology alone has no power: every alignment family in
use is a homeomorphism, and on finite data a manufactured alignment can be
realized by a diffeomorphism isotopic to the identity
(Appendix~\ref{app:topo-blind}). Metric persistence does constrain the map.
Interchange interventions impose a \emph{stretch budget}: the map's distortion
times the sensitivity of the downstream computation must exceed the output
margin divided by the network's own separation of the values of $V$, which is
the first bar of a chromatic persistence diagram
(Appendix~\ref{app:topo-stretch}). Persistence diagrams of the aligned cloud
bound distortion only weakly (Appendix~\ref{app:topo-interleave}). Topology
beyond connected components enters in two narrower places: whether
counterfactual states stay on the network's natural manifold
(Appendix~\ref{app:topo-real}), and variables whose values carry topology of
their own (Appendix~\ref{app:topo-struct}).
Appendix~\ref{app:topo-codes} relates the stretch budget to the
prior-dependent terms of Stage~1 code lengths, and
Appendix~\ref{app:topo-protocol} turns the results into diagnostics, two
hypotheses and a proof target.

\subsection{Setting}
\label{app:topo-setup}

Fix a locus $\ell$. Write $h(x)\in\mathbb R^d$ for the state on input $x$,
$H=\{h(x_i)\}_{i\le n}$ for the finite cloud of states, and
$g:\mathbb R^d\to\mathbb R^K$ for the network's fixed computation from the
locus to the logits. An alignment map $\phi$ is invertible, with
$\phi(h)=(z_V,z_\perp)$ and $z_V\in\mathbb R^{d_V}$; the rotations of
distributed alignment search and the reversible residual networks of
\cite{sutter2025} are both of this kind. The intervened state for base $b$ and
source $s$ is
\begin{equation}
h'(b,s)=\phi^{-1}\bigl(z_V(s),\,z_\perp(b)\bigr).
\label{eq:topo-interv}
\end{equation}
Norms are Euclidean. The Lipschitz constant of a map $f$ on a set $S$ and the
distortion of the alignment map are
\[
\Lip_S(f)=\sup_{x\neq y\in S}\frac{\|f(x)-f(y)\|}{\|x-y\|},\qquad
\Dist(\phi)=\Lip(\phi)\Lip(\phi^{-1}),
\]
and the sample distortion is
$\Dist_H(\phi)=\Lip_H(\phi)\Lip_{\phi(H)}(\phi^{-1})$. $\VR_t(S)$ is the
Vietoris--Rips complex of simplices of diameter at most $t$, and
$\Dgm_j(S)$ its persistence diagram in degree $j$. Colouring each state by
$V_A(x)$ gives classes $H_v$; the inclusion
$\bigsqcup_v\VR_t(H_v)\hookrightarrow\VR_t(H)$ defines kernel, image and
cokernel persistence modules, the objects of chromatic persistence
\cite{cultrera2025chromatic}.

\paragraph{Margins.} For a base $b$ and sources $s,s'$ whose counterfactual
labels $y\neq y'$ differ, call both interventions correct with margin $\mu$ if
in each the target logit exceeds every other logit by at least $\mu$. The
logit difference $\delta=g(h'(b,s))-g(h'(b,s'))$ then satisfies
$\delta_y-\delta_{y'}\ge2\mu$, hence $\|\delta\|\ge\gamma:=\sqrt2\,\mu$. A
prediction that puts probability at least $1-\eta$ on the target has
$\mu\ge\log\bigl((1-\eta)/\eta\bigr)$. A prediction that reproduces the
randomized-response target law of Section~\ref{sec:stage1} exactly, with $K$
labels and noise rate $\eps_{\rm e}$, has
$\mu=\log\bigl(1+K(1-\eps_{\rm e})/\eps_{\rm e}\bigr)$.

\subsection{Topology alone cannot separate genuine from manufactured alignments}
\label{app:topo-blind}

\LCalc{} (i) A homeomorphism preserves every homotopy invariant. If $M$ is the
support of the states, $\phi(M)$ and $M$ have the same Betti numbers however
$\phi$ was found. Topology changes only in maps that are not homeomorphisms,
such as the network's own ReLU layers \cite{naitzat2020}.

(ii) \emph{Finite realizability.} Let $d\ge2$ and $1\le d_V<d$. Suppose that
for every label $y$ the set
$\Omega_y=\{h:\ g(h)_y>\max_{j\neq y}g(h)_j+\mu\}$ is nonempty, a margin
version of the output-surjectivity assumed by \cite{sutter2025}; it is open
because $g$ is continuous. Suppose the states of distinct inputs are distinct
(input-injectivity). Then for any algorithm $A$ and any finite set
$\mathcal P$ of pairs $(b,s)$ with $s\neq b$, there is a compactly supported
$C^\infty$ diffeomorphism of $\mathbb R^d$, isotopic to the identity, whose
interchange interventions are correct with margin $\mu$ on every pair in
$\mathcal P$.

\emph{Proof.} Choose distinct $u_x\in\mathbb R^{d_V}$ and distinct
$w_x\in\mathbb R^{d-d_V}$ for the inputs. The data codes $(u_x,w_x)$ and the
recombined codes $(u_s,w_b)$, $(b,s)\in\mathcal P$, are pairwise distinct. For
each pair pick $q_{b,s}\in\Omega_{A(b;V\leftarrow V_A(s))}$, distinct from one
another and from the states, which is possible because these sets are open
and nonempty. The prescription $h(x)\mapsto(u_x,w_x)$,
$q_{b,s}\mapsto(u_s,w_b)$ is a bijection between finite sets of distinct
points. For $d\ge2$ the space of configurations of finitely many distinct
labelled points is connected, so by the isotopy extension theorem
\cite{hirsch1976} the bijection extends to a diffeomorphism of the required
kind, and $h'(b,s)=q_{b,s}$. $\square$

Coupling-based invertible networks approximate such diffeomorphisms
\cite{teshima2020}; an approximation error at $q_{b,s}$ smaller than its
distance to the boundary of $\Omega_y$ leaves the intervention correct.

\paragraph{Consequences.} The non-linear representation dilemma survives
restriction to maps that deform smoothly to the identity. Because the codes
are arbitrary, $\phi(H)$ can be any configuration of $n$ distinct points, and
$\Dgm_j(\phi(H))$ any diagram that $n$ points realize. A test that asks only
whether a map preserves topological features, or that reads the topology of
the aligned cloud alone, has no power. Evidence can come only from scales.

\subsection{The stretch budget}
\label{app:topo-stretch}

\begin{keybox}{Stretch budget (internal calculation).}
Let $b,s,s'$ be as in Appendix~\ref{app:topo-setup}, with both interventions
correct with margin $\mu$, and let
$\Lambda_{b,s,s'}=\sup_{\vartheta\in[0,1]}\bigl\|Dg\bigl((1-\vartheta)h'(b,s)+\vartheta h'(b,s')\bigr)\bigr\|$
be the sensitivity of $g$ along the segment between the two intervened
states. Then
\begin{equation}
\begin{aligned}
\gamma&\le\Lambda_{b,s,s'}\,\|h'(b,s)-h'(b,s')\|
\le\Lambda_{b,s,s'}\Lip(\phi^{-1})\,\|z_V(s)-z_V(s')\|\\
&\le\Lambda_{b,s,s'}\Dist(\phi)\,\|h(s)-h(s')\| .
\end{aligned}
\label{eq:topo-chain}
\end{equation}
Write
\[
\tau_{\rm f}=\frac{\|z_V(s)-z_V(s')\|}{\|h(s)-h(s')\|},\qquad
\tau_{\rm i}=\frac{\|h'(b,s)-h'(b,s')\|}{\|z_V(s)-z_V(s')\|}
\]
for the forward and inverse stretch, and $\Lambda_{\rm nat}$ for a reference
sensitivity of $g$ on the natural states, such as the median of
$\|Dg(h(x))\|$. The chain is then a budget,
\begin{equation}
\underbrace{\log\frac{\gamma}{\Lambda_{\rm nat}\,\|h(s)-h(s')\|}}_{\text{set by network and targets}}
\ \le\ \log\tau_{\rm f}+\log\tau_{\rm i}+\log\frac{\Lambda_{b,s,s'}}{\Lambda_{\rm nat}} .
\label{eq:topo-budget}
\end{equation}
Minimizing over pairs gives
\begin{equation}
\Dist(\phi)\,\Lambda\ \ge\ \frac{\gamma}{\sep_V},\qquad
\sep_V=\min\bigl\{\|h(s)-h(s')\|:\ V_A(s)\neq V_A(s')\bigr\},
\label{eq:topo-sep}
\end{equation}
where $\Lambda$ is the largest segment sensitivity and the minimum runs over
pairs whose counterfactual labels differ on some base and whose interventions
are correct.
\end{keybox}

\noindent\emph{Proof.} The first inequality is the mean value inequality
together with $\|\delta\|\ge\gamma$. The second applies $\phi^{-1}$ to two
codes that differ only in $z_V$. The third uses
$\|z_V(s)-z_V(s')\|\le\|\phi(h(s))-\phi(h(s'))\|$. The budget is the chain
rewritten with
$\|h'(b,s)-h'(b,s')\|=\tau_{\rm i}\tau_{\rm f}\|h(s)-h(s')\|$. $\square$

\paragraph{The separation is a persistence quantity.} \LCalc{} In degree zero,
the kernel of $\bigsqcup_v\VR_t(H_v)\hookrightarrow\VR_t(H)$ has dimension
$\sum_v\beta_0(\VR_t(H_v))-\beta_0(\VR_t(H))$, because every component of
$\VR_t(H)$ contains a monochromatic component. For $t<\sep_V$ every path in
$\VR_t(H)$ is monochromatic, so the kernel vanishes; at $t=\sep_V$ the closest
bichromatic edge joins two classes. When all pairs of distinct values
qualify, $\sep_V$ is therefore the first birth in the degree-zero kernel
diagram, and the rest of that barcode records the scales at which the classes
mingle. \LNum{} The identity holds exactly in 20 random three-colour clouds
(Appendix~\ref{app:topo-checks}).

\paragraph{Reading the budget.} Define the separation ratio
$\rho_V=\Lambda_{\rm nat}\sep_V/\gamma$. If $\rho_V\ge1$, the network already
separates $V$ at the scale its downstream computation reads, and a rotation
satisfies \eqref{eq:topo-sep}. If $\rho_V<1$, any alignment that succeeds must
supply $\log_2(1/\rho_V)$ bits of stretch: from the forward map, from the
inverse map, or by sending intervened states to regions where $g$ is more
sensitive than on the natural states. All three terms of
\eqref{eq:topo-budget} can be measured pair by pair. An alignment can lower
its budget only by giving up the closest bichromatic pairs, whose
interventions then count as failures; when interchange-intervention accuracy
is below one, $\sep_V$ is computed over the sources whose interventions
succeed.

\paragraph{Geometry that is blind to $V$.} \LCalc{} Suppose the states are $n$
independent draws from a density $f$ on a $k$-dimensional region and the
labels are independent of position, with class probabilities $\pi_v$. For
small $r$ the number of bichromatic pairs within distance $r$ is then
approximately Poisson \cite{jammalamadaka1986} with mean
$\binom n2\bigl(1-\sum_v\pi_v^2\bigr)\,\omega_kr^k\!\int\!f^2$, where
$\omega_k$ is the volume of the unit ball. The median of $\sep_V$ is therefore
of order $n^{-2/k}$, and \eqref{eq:topo-sep} requires
\begin{equation}
\Dist(\phi)\,\Lambda\ \gtrsim\ \gamma\,n^{2/k},
\label{eq:topo-blind}
\end{equation}
whereas a network that computes $V$ with margin $\mu_0$ keeps
$\sep_V\ge\mu_0$ for every $n$. \LNum{} On the flat torus with uniform
density, fitted slopes of $\log\sep_V$ against $\log n$ for
$n=2^8,\dots,2^{15}$ lie within $0.012$ of $-2/k$ for $k=2,3,4$, and the
median of $\sep_V$ lies within 10\% of the Poisson prediction
(Appendix~\ref{app:topo-checks}).

\paragraph{Discrete inputs.} With finitely many token inputs, $\sep_V$ stops
shrinking at the distance between inputs that differ only in $V$, and the
budget depends on how the network treats $V$. \LHyp{} For computed
variables---the equality bits of hierarchical equality, or any variable at a
locus before it is computed---random features interleave the classes and
$\rho_V\ll1$. For copied variables such as the indirect-object name, random
networks preserve token identity, $\rho_V$ may be close to one, and
\eqref{eq:topo-sep} does not bite. This is consistent with the perfect
alignments that \cite{sutter2025} obtained on that task in randomly
initialized models.

\subsection{Persistence of the aligned cloud bounds distortion only weakly}
\label{app:topo-interleave}

\LCalc{} Let $L=\Lip_H(\phi)$ and $L'=\Lip_{\phi(H)}(\phi^{-1})$. The vertex
maps give simplicial maps $\VR_t(H)\to\VR_{Lt}(\phi(H))\to\VR_{LL't}(H)$ whose
composite is the inclusion, and likewise in the other order. Rescaling
$\phi(H)$ by $\sqrt{L'/L}$ turns this into a multiplicative interleaving with
factor $\sqrt{\Dist_H(\phi)}$, which on a logarithmic scale is an additive
interleaving with shift $\tfrac12\log\Dist_H(\phi)$. By the isometry theorem
\cite{chazal2016}, for every degree $j$,
\begin{equation}
d_B\Bigl(\log\Dgm_j(H),\ \log\Dgm_j\bigl(\phi(H)\bigr)+\tfrac12\log(L'/L)\Bigr)
\ \le\ \tfrac12\log\Dist_H(\phi),
\label{eq:topo-interleave}
\end{equation}
where bars in degree zero, all born at zero, are compared by their deaths.
Since $\phi$ preserves the colouring, the same bound holds for the kernel
modules of Appendix~\ref{app:topo-stretch}. \LNum{} The bound held in 60
trials with four clouds and coupling, linear and radial maps, with
$\Dist_H(\phi)$ from 1.04 to 235. It was attained up to rounding in degree
zero; the median ratio of the two sides was 0.72 in degree zero and 0.59 in
degree one (Appendix~\ref{app:topo-checks}).

The bound runs one way only. A large diagram distance proves a large
distortion, but $\Dist_H(\phi)$ itself costs $O(n^2)$ distance evaluations and
is never smaller than the bound implies. On the map side, persistence adds
robustness to single extreme pairs and nothing more. The quantity to report
is the sample distortion, together with the segment sensitivities of
\eqref{eq:topo-budget}.

\subsection{Realizability of counterfactual states}
\label{app:topo-real}

Causal abstraction constrains outputs only. The stronger hypothesis that the
representation factorizes through the algorithm's variables also constrains
where intervened states lie, and here topology beyond connected components
becomes informative.

\begin{keybox}{Realizability bound (internal calculation).}
Suppose there are a map $\psi$, codes $c$ on the values of $V$ and $c'$ on the
remaining variables $W$, with $\phi(\psi(v,w))=(c(v),c'(w))$ and
$\|h(x)-\psi(V_A(x),W(x))\|\le\eps$ for all inputs. Suppose the input set
contains, for each pair $(b,s)$, a recombined input $x^*$ with
$V_A(x^*)=V_A(s)$ and $W(x^*)=W(b)$. Then
\begin{equation}
\|h'(b,s)-h(x^*)\|\ \le\ \bigl(\sqrt2\,\Dist(\phi)+1\bigr)\,\eps .
\label{eq:topo-real}
\end{equation}
The intervened cloud $I$ therefore lies within
$(\sqrt2\,\Dist(\phi)+1)\eps$ of $H$ in Hausdorff distance, and by the
stability theorem for Rips complexes \cite{chazal2014}
\begin{equation}
d_B\bigl(\Dgm_j(H),\Dgm_j(H\cup I)\bigr)\ \le\ 2\bigl(\sqrt2\,\Dist(\phi)+1\bigr)\,\eps
\qquad\text{for all }j .
\label{eq:topo-real-dgm}
\end{equation}
\end{keybox}

\noindent\emph{Proof.} Both $\|z_V(s)-c(V_A(s))\|$ and
$\|z_\perp(b)-c'(W(b))\|$ are at most $\Lip(\phi)\eps$, so the code
$(z_V(s),z_\perp(b))$ lies within $\sqrt2\Lip(\phi)\eps$ of
$(c(V_A(s)),c'(W(b)))=\phi(\psi(V_A(s),W(b)))$. Applying $\phi^{-1}$
multiplies this by at most $\Lip(\phi^{-1})$, and $\psi(V_A(s),W(b))$ lies
within $\eps$ of $h(x^*)$. The diagram bound follows because the
Gromov--Hausdorff distance between $H$ and $H\cup I$ is at most their
Hausdorff distance. $\square$

\LNum{} In 60 trials with a smooth bi-Lipschitz factorization, seven cyclic
values and $\eps$ from 0.01 to 0.2, the largest observed error was 0.59 of the
bound in \eqref{eq:topo-real}, and \eqref{eq:topo-real-dgm} held in degrees
zero and one (Appendix~\ref{app:topo-checks}).

For manufactured alignments nothing ties $I$ to $H$: in the construction of
Appendix~\ref{app:topo-blind} the intervened states $q_{b,s}$ sit anywhere in
the label regions $\Omega_y$. Three statistics follow.
\begin{enumerate}[leftmargin=*]
\item The representation-level counterfactual error $\|h'(b,s)-h(x^*)\|$
where recombined inputs exist, normalized by the nearest-neighbour spacing of
$H$.
\item The one-sided Hausdorff distance from $I$ to $H$ where they do not.
\item The persistence of the inclusion $H\hookrightarrow H\cup I$, a
1-Lipschitz map of the kind studied in \cite{gozzi2025}. Its cokernel bars in
degree zero are components made of intervened states only; its kernel bars in
degree one are loops of the natural manifold that intervened states fill in.
\end{enumerate}
The last is the genuinely topological signal. Interventions that cut across a
curved representation instead of following it, as linear steering does
\cite{bhalla2026}, produce it. A failure indicts factorization, which is
stronger than causal abstraction: subspace patching can leave the manifold
even when the network computes $V$ \cite{makelov2024}.

\subsection{Variables whose values carry topology}
\label{app:topo-struct}

Some variables take values in a space with its own geometry, such as the
circular day-of-week and month features of \cite{engels2025} or residues in
modular arithmetic. A stronger, geometric abstraction can then be tested: the
network's class space should reproduce the value space up to bounded
distortion.

\LCalc{} With the recombined inputs of Appendix~\ref{app:topo-real}, define
the class dissimilarity $d^V(v,v')$ as the median over $w$ of
$\|h(x_{v,w})-h(x_{v',w})\|$, the distances between inputs that differ only
in $V$. Under $\phi$ each such distance is multiplied by a factor in
$[1/\Lip(\phi^{-1}),\Lip(\phi)]$, and so is any order statistic of them. The
argument of \eqref{eq:topo-interleave} needs only symmetry, so it applies to
Rips filtrations of these dissimilarities. Two consequences follow.
\begin{enumerate}[leftmargin=*]
\item \emph{Survival.} If $a\,d_{\mathcal V}\le d^V\le a\,\Dist_V\,d_{\mathcal V}$
for some $a>0$, where $d_{\mathcal V}$ is the metric of the value space, then
a degree-one bar of log-length $\xi>\log\Dist_V$ in the value space appears
in the network's class space with log-length at least $\xi-\log\Dist_V$. For $p$
evenly spaced values on a circle the bar has death-to-birth ratio
$\sin(\pi\lceil p/3\rceil/p)/\sin(\pi/p)$ \cite{adamaszek2017}: 2.25 for
$p=7$, 3.35 for $p=12$ and 6.63 for $p=24$. Short cycles therefore survive
only small distortions. \LNum{} Rips computations reproduce this ratio for
$p=4,\dots,24$ (Appendix~\ref{app:topo-checks}).
\item \emph{Detection.} If the network's class space has no degree-one bar of
log-length above $\eta$, while the class space computed in full
$\phi$-coordinates has one of log-length $\xi$, then
$\log\Dist(\phi)\ge\xi-\eta$.
\end{enumerate}
Two cautions apply. The projection to $z_V$ is 1-Lipschitz but can create
loops, as a helix projects to a circle; class spaces are therefore computed in
$h$ and in full $\phi$-coordinates, and a loop that appears only after
projection is a warning sign. And the persistence of a class space does not
depend on which value is which. A test of whether the network's class space
matches the value space must use the correspondence, for example the
distortion of $v\mapsto v$ between $d^V$ and $d_{\mathcal V}$ or the
cross-barcode of \cite{barannikov2022}; permuting the correspondence then
gives an exact null distribution.

\subsection{Relation to the code lengths of Stage 1}
\label{app:topo-codes}

In the Occam regime of Section~\ref{sec:stage1}, configurations that reproduce
the same target law differ only through logarithmic and prior-dependent
terms. The stretch budget says where a manufactured configuration pays: in
the prior mass of maps with high distortion. Below, $\Pi$ is the prior on the
parameters of the alignment map.

\begin{keybox}{Prior mass of high distortion (internal calculation).}
Let $\phi$ be an additive-coupling network \cite{dinh2015,gomez2017} with $B$
blocks
\[
x\ \mapsto\ O_b\bigl(x_1,\ x_2+\Theta_{b2}\,\mathrm{ReLU}(\Theta_{b1}x_1+\nu_b)\bigr),
\qquad \Theta_{bi}\in\mathbb R^{d_{bi}\times d'_{bi}},
\]
with fixed permutation matrices $O_b$ and independent $N(0,\sigma^2)$ priors
on the entries of the matrices $\Theta_{bi}$. Each block and
its inverse are $(1+\|\Theta_{b1}\|\|\Theta_{b2}\|)$-Lipschitz, so
\[
\Dist(\phi)\ \le\ \prod_b\bigl(1+\|\Theta_{b1}\|\|\Theta_{b2}\|\bigr)^2 .
\]
Let
$s_{\max}=\max_{b,i}\sigma\bigl(\sqrt{d_{bi}}+\sqrt{d'_{bi}}\bigr)$ and
$u=(\Dist^*)^{1/(2B)}-1$. If $u\ge4s_{\max}^2$, then
\begin{equation}
-\log\Pi\bigl(\Dist(\phi)\ge \Dist^*\bigr)\ \ge\ \frac{u}{8\sigma^2}-\log(2B).
\label{eq:topo-prior}
\end{equation}
\end{keybox}

\noindent\emph{Proof.} If $\Dist(\phi)\ge \Dist^*$, some block has
$\|\Theta_{b1}\|\|\Theta_{b2}\|\ge u$, so one of its two matrices has norm at
least $\sqrt u\ge2\sigma(\sqrt{d_{bi}}+\sqrt{d'_{bi}})$. Gaussian
concentration of the operator norm,
$\Pr\bigl(\|\Theta\|\ge\sigma(\sqrt{d}+\sqrt{d'})+a\bigr)\le e^{-a^2/(2\sigma^2)}$
\cite{vershynin2018}, with $a\ge\sqrt u/2$, and a union bound over the $2B$
matrices give the claim. $\square$

\paragraph{Code lengths.} \LCalc{} (sketch). Let $\mathcal A$ be the set of
maps whose interventions are correct with margin $\mu$ on all $n$ pairs. If
$g$ is globally $\Lambda_g$-Lipschitz, \eqref{eq:topo-sep} gives
$\mathcal A\subseteq\{\Dist(\phi)\ge \Dist^*\}$ with $\Dist^*=\gamma/(\Lambda_g\sep_V)$.
Assume that maps outside $\mathcal A$ pay at least $\Delta_n$ nats more than
the best fit $nL^*$. Then
\[
F_n\ \ge\ nL^*-\log\bigl[\Pi(\mathcal A)+e^{-\Delta_n}\bigr],
\]
so when $\Delta_n\ge-\log\Pi(\mathcal A)$ the code length exceeds the best
fit by at least $-\log\Pi(\Dist(\phi)\ge \Dist^*)-\log2$, which
\eqref{eq:topo-prior} bounds from below. Three readings follow.
\begin{enumerate}[leftmargin=*]
\item The prior-dependent $O_p(1)$ terms of \eqref{eq:gain} are where the
stretch is paid, at least about $(\Dist^*)^{1/(2B)}/(8\sigma^2)$ nats.
\item Depth makes distortion cheap: each block takes a root of $\Dist^*$. This is
a quantitative reading of the observation that more complex maps align even
untrained networks \cite{sutter2025}. For deep maps the prior term will rarely
decide a comparison, which is why the diagnostics below report distortion
directly.
\item Under geometry blind to $V$ with continuous inputs,
\eqref{eq:topo-blind} gives $\Dist^*\propto n^{2/k}$ and a term growing like
$n^{1/(kB)}$. In the population limit no map attains the margin, and the
configuration leaves the Occam regime for the falsification regime.
\end{enumerate}
Whether distortion is linked to the learning coefficient of the map family
remains open.

\subsection{Diagnostics, hypotheses and proof target}
\label{app:topo-protocol}

\LTest{} Three diagnostics accompany every alignment in the decision
experiment. Together they cost far less than training the alignments.
\begin{description}[style=nextline,leftmargin=1.5em]
\item[D1, network side.] The $V$-coloured kernel barcode in degree zero at each
locus, computed from a minimum spanning tree with union--find; $\sep_V$, its
quantiles and $\rho_V$; and a null distribution from shuffling $V$-labels
across inputs, which is exact under exchangeability and gives a p-value or,
through a calibrator, an e-value. Rejecting this null shows that $V$ is
represented metrically, not that an abstraction is valid, so the statistic is
compared with the untrained-network control on the same inputs.
\item[D2, map side.] The three terms of \eqref{eq:topo-budget} for every pair
used in evaluation: forward and inverse stretch exactly, and segment
sensitivity by power iteration on Jacobian--vector products at several points
per segment. Sampling underestimates the supremum, so a budget that the
measured terms do not account for points to sensitivity spikes between the
sample points. The sample distortion $\Dist_H(\phi)$ is reported alongside.
\item[D3, realizability.] The counterfactual error of \eqref{eq:topo-real}
where recombined inputs exist, the one-sided Hausdorff distance otherwise,
and persistence of $H\hookrightarrow H\cup I$ in degrees zero and one on
subsamples of at most a few thousand states. For variables whose values carry
topology, the class-space comparison of Appendix~\ref{app:topo-struct} is
added.
\end{description}
D1 is observational. The benchmark of \cite{meloux2026} found that only causal
metrics with faithfulness tests on unmapped variables discriminated reliably,
so D1 alone is not expected to discriminate. D2 and D3 use the interchange
constraint; they enter the decision experiment as additional rivals and as
candidate components of the complete protocol of H-A4.

\begin{hypbox}{H-A8 (stretch budgets separate computed variables).}
On configurations with valid explanations in the decision experiment, the
network-side term of \eqref{eq:topo-budget} is nonpositive on most evaluation
pairs. On configurations with manufactured alignments of computed variables---the
equality bits of hierarchical equality on untrained networks, and alignments
at loci before $V$ is computed---it is positive, and D2 attributes it to
stretch or to off-manifold sensitivity, including for alignments that pass
i.i.d.\ held-out IIA. For copied variables such as the indirect-object name,
budgets are not expected to separate.
\end{hypbox}

\begin{hypbox}{H-A9 (counterfactual states).}
Valid alignments on ground-truth models keep the counterfactual error
\eqref{eq:topo-real} within a small multiple of the nearest-neighbour spacing
of $H$ and leave the persistence of $H\hookrightarrow H\cup I$ nearly trivial.
Manufactured alignments that pass distribution-shifted held-out IIA,
including those for copied variables, place intervened states off the natural
manifold.
\end{hypbox}

\begin{keybox}{Proof target T7: metric certificates.}
\begin{description}[style=nextline,leftmargin=1.2em]
\item[T7a, finite realizability by diffeomorphisms.] Status: proved
(Appendix~\ref{app:topo-blind}).
\item[T7b, stretch budget and its persistence form.] Status: proved
(Appendix~\ref{app:topo-stretch}); the rate for geometry blind to $V$ rests on
a standard Poisson approximation \cite{jammalamadaka1986} and is checked
numerically.
\item[T7c, interleaving and realizability bounds.] Status: proved from
standard stability results (Appendices~\ref{app:topo-interleave}
and~\ref{app:topo-real}).
\item[T7d, prior mass of distortion.] Status: proved for additive-coupling
networks with Gaussian priors; the consequence for code lengths needs the fit
assumption on $\Delta_n$ to be established for a concrete family.
\item[T7e, distortion and learning coefficients.] Whether the local learning
coefficient of an alignment family grows with the distortion of the maps that
fit the targets. Status: open.
\end{description}
\end{keybox}

\subsection{Numerical checks}
\label{app:topo-checks}

\LNum{} The checks concern the stated models only. Rips persistence and
bottleneck distances use the \texttt{gudhi} library, version~3.13.

\begin{center}\small
\begin{tabular}{@{}P{5.9cm}P{9.6cm}@{}}
\toprule
Claim and script & Result \\
\midrule
Interleaving bound \eqref{eq:topo-interleave} in degrees zero and one; first
kernel birth equals $\sep_V$. \newline\texttt{topo\_interleave\_check.py}
& Held in 60 of 60 trials with $\Dist_H(\phi)$ from 1.04 to 235; largest ratio
of the two sides 1.000 in degree zero and 0.996 in degree one. Kernel identity
exact in 20 of 20 clouds. \\ \addlinespace
Poisson law for $\sep_V$ under geometry blind to $V$.
\newline\texttt{topo\_bichromatic\_check.py}
& Fitted slopes $-0.991$, $-1.000$ ($k=2$), $-0.655$, $-0.669$ ($k=3$),
$-0.502$, $-0.509$ ($k=4$), for two and three equiprobable classes; median
over prediction between 0.91 and 1.09 for all $n$ (200 replicates each). \\
\addlinespace
Realizability bounds \eqref{eq:topo-real} and \eqref{eq:topo-real-dgm}.
\newline\texttt{topo\_realizability\_check.py}
& Held in 60 of 60 trials with $\Dist(\phi)\le(1+\beta)/(1-\beta)$,
$\beta$ from 0.1 to 0.9; largest error over bound 0.59, median 0.27. \\
\addlinespace
Death-to-birth ratio of the degree-one bar for $p$ evenly spaced values.
\newline\texttt{topo\_circle\_check.py}
& Matches $\sin(\pi\lceil p/3\rceil/p)/\sin(\pi/p)$ to four decimals for
$p=4,5,6,7,8,10,12,24$. \\
\bottomrule
\end{tabular}
\end{center}

\section{Notation and evidence labels}

\begin{longtable}{@{}P{3.2cm}P{11.7cm}@{}}
\toprule
Symbol or term & Meaning in this plan \\
\midrule
\endhead
$\lambda$, $m$ & Learning coefficient (real log canonical threshold) and
multiplicity, global or local to a region or solution. \\ \addlinespace
$F_n$, $F_U$ & Bayes free energy, the length of the mixture code; its
restriction to a region $U$. \\ \addlinespace
$L^A(n)$ & Prequential code length of learner $A$ on $n$ labels. \\ \addlinespace
$\DL(n)$, $G(n)$ & Stage~1: prequential code length of fixed counterfactual targets under a network with a learned alignment map; gain over a control on the same targets, in bits. \\ \addlinespace
$\Delta\ell$, $\Delta\lambda$ & Faithfulness gap and complexity gap
between a control and a hypothesis. \\ \addlinespace
Controls & Untrained network, wrong locus or other map family on the same targets; a label-only code without access to the algorithm. \\ \addlinespace
$r$, $\kappa$ & Growth factor of prequential blocks, which multiplies logarithmic coefficients by $(r-1)/\ln r$; a learner's generalization coefficient. \\ \addlinespace
EDL, $\mathrm{AC}_A$ & Excess description length, the prequential code length minus $n$ times the final test loss; the algorithmic predictive cost of learner $A$ relative to the Bayes predictor. \\ \addlinespace
$D$, $k^*$ & Depth, the log-odds against the capable region under the local posterior; the number of demonstrations needed to flip it. \\ \addlinespace
$E$ & An e-value: a nonnegative statistic with expectation at most one under
the null; $\log_2E$ is evidence in bits. \\ \addlinespace
$\chi_\phi$ & Susceptibility of an observable $\phi$ to a shift of the data distribution. \\ \addlinespace
IIA & Interchange-intervention accuracy. \\ \addlinespace
$\Dist(\phi)$, $\Dist_H(\phi)$ & Distortion $\Lip(\phi)\Lip(\phi^{-1})$ of an
alignment map, globally and on the sample of states
(Appendix~\ref{app:topo}). \\ \addlinespace
$\sep_V$, $\rho_V$ & Smallest distance between states of inputs with
different values of $V$, the first bar of the $V$-coloured kernel diagram;
separation ratio $\Lambda_{\rm nat}\sep_V/\gamma$. \\ \addlinespace
$\gamma$, $\Lambda$ & Output margin in logit space; sensitivity of the
network's computation downstream of the locus. \\ \addlinespace
$\chi_{\rm exit}$ & Scaling exponent of the exit cross-section of a
flat-set tube. \\ \addlinespace
$\hat T$, $\hat\lambda$ & Thermometer estimates from the first two
moments of the excess loss. \\ \addlinespace
Evidence labels & Known result, internal calculation, numerical check,
hypothesis, proposed test and assessment, as defined in
Section~\ref{sec:purpose}; a dagger marks unverified bibliographic
details. \\
\bottomrule
\end{longtable}

\section{Change log}

\paragraph{Version 7.2 (12 October 2026): topological and metric certificates.}
\begin{enumerate}[leftmargin=*]
\item Appendix~\ref{app:topo} asks whether persistent homology can test
alignment maps. Topology alone cannot: every alignment family in use is a
homeomorphism, and manufactured alignments are realizable by diffeomorphisms
isotopic to the identity.
\item Interchange interventions impose a stretch budget relating the map's
distortion, the downstream sensitivity and the network's separation of the
values of $V$, which is the first bar of a chromatic persistence diagram.
\item A realizability bound for counterfactual states, a class-space test for
variables whose values carry topology, and a prior-mass bound linking
distortion to Stage~1 code lengths are added, with diagnostics D1--D3,
hypotheses H-A8 and H-A9, and proof target T7.
\end{enumerate}

\paragraph{Version 7.1 (24 September 2026): second review.}
\begin{enumerate}[leftmargin=*]
\item Every gate has a stop rule: failing lines stop as paper projects,
further work needs a bounded feasibility budget and a new hypothesis with
preliminary evidence, and Stage~2 needs its own pilot. Promises that every
outcome is publishable are removed.
\item Project~B separates recoverability, the primary outcome, from the
training history of organisms, a secondary analysis.
\item Products of e-values are used only with proven conditional validity;
otherwise evidence is combined by prespecified weighted averages.
\item The batching remainder is $o(\log n)$ in general and $O(1)$ under a
stronger assumption; the Occam regime requires equal conditional target
laws; derivatives are taken along a continuous population-loss tilt.
\item Stage~2 fixes its unit of observation and freezes its intervention
protocol before scoring; H-A6 is softened.
\end{enumerate}

\paragraph{Version 7.0 (24 September 2026): response to an external review.}
\begin{enumerate}[leftmargin=*]
\item Project~A becomes two-staged. Stage~1 measures alignment sample
efficiency on fixed targets, with controls whose information access is
specified; Stage~2 codes the network's own interventional behaviour with
competing explanations on shared data (Section~\ref{sec:projA}).
\item Close prior work on prequential causal description length is
incorporated, and distribution-shifted held-out accuracy and complete
intervention protocols become primary rivals.
\item A theory strand is added: the batching factor, derived and checked;
learner coefficients; prior terms; and certified tests
(Section~\ref{sec:atests}; Appendix~\ref{app:batch}).
\item Corrections: EDL subtracts test loss; perfect accuracy does not imply
equal log loss; smaller maps do not imply smaller learning coefficients;
universal inference needs certified bounds; sequential e-values need fresh
data; and adding demonstrations differs from replacing data, with the flip
now computed exactly by reweighting.
\item Project~B is narrowed to one geometric statistic, validated on exact
masses before use, with recoverability as the operational target
(Section~\ref{sec:projB}).
\item The plan opens with a six-week decision experiment, and publication
probabilities are replaced by scenarios and readiness indicators
(Sections~\ref{sec:schedule} and~\ref{sec:outlook}).
\end{enumerate}

\paragraph{Version 6.0 (from version 5).}

\begin{enumerate}[leftmargin=*]
\item \textbf{Organizing principle.} The question changes from whether SLT
invariants control training dynamics to how much evidence, in bits,
supports an interpretability or capability claim. SLT supplies the
complexity term (Sections~\ref{sec:purpose} and~\ref{sec:primer}).
\item \textbf{Flagship.} New Project~A on the interventional
description-length gain, addressing the non-linear representation
dilemma, with a decomposition into faithfulness and complexity gaps and
ground-truth tests (Section~\ref{sec:projA}; Appendix~\ref{app:codes}).
\item \textbf{Safety project.} New Project~B on suppressed versus absent
capabilities, with an exact split of excess description length and local
free-energy measures of suppression depth. The emergent-misalignment
study becomes a case (Section~\ref{sec:projB}).
\item \textbf{Dynamics.} Closed in a short note: scoped impossibility,
corrected flat-set law and discrete-chain corrections
(Section~\ref{sec:p0}; Appendix~\ref{app:p0}).
\item \textbf{Instruments.} The thermometer and identifiability checks
are retained as calibration tools, with their limits stated
(Appendix~\ref{app:calib}).
\item \textbf{Literature.} Incorporates the September 2026 sweep on
causal abstraction, description-length interpretability, elicitation,
sandbagging audits, stochastic-gradient MCMC at scale and emergent
misalignment.
\item \textbf{Planning.} Effort, compute, contacts, gates, descoping and
publication probabilities are made explicit
(Sections~\ref{sec:schedule}--\ref{sec:risks}).
\item \textbf{Conventions.} Evidence labels now include ``hypothesis'' and
``assessment''.
\end{enumerate}

\paragraph{Revision 6.1 (23 September 2026): full-time commitment.}
\begin{enumerate}[leftmargin=*]
\item The effort assumption changes from half time to full time; compute
is available and is no longer a planning constraint
(Section~\ref{sec:purpose}).
\item Project~A splits into A1, a method paper aimed at ICML 2027 with
NeurIPS 2027 as a second chance, and A2, a language-model-scale safety
application with proofs. New hypothesis H-A5 and proof targets T1--T4
(Section~\ref{sec:a1a2}).
\item Project~B starts in month 3, runs up to eight billion parameters, and
aims at NeurIPS 2027 as a stretch and ICLR 2028 as the base
(Section~\ref{sec:projB}).
\item Gates move earlier. The descoping ladder is replaced by reserve
release rules and a two-project concurrency rule
(Section~\ref{sec:schedule}).
\item The publication outlook and risks are updated
(Sections~\ref{sec:outlook} and~\ref{sec:risks}).
\end{enumerate}

\paragraph{Revision 6.2 (23 September 2026): corrections after internal review.}
\begin{enumerate}[leftmargin=*]
\item The fixed-error-rate likelihood behind the gain decomposition was
not analytic in the parameters. It is replaced by randomized-response
targets, which make the model realizable with interior optima
(Appendix~\ref{app:codes}; proof target T1).
\item Two hypotheses were true by construction, because optimizer effects
on EDL lie in the algorithmic term by definition. H-B1 and H-B4 now make
testable quantitative claims, and the EDL identity is stated as exact only
under a Bayesian data model (Section~\ref{sec:projB}).
\item Held-out IIA, the generalization criterion suggested by Sutter et
al., becomes the decisive rival for Project~A, with an argument for why
the prequential code should beat it (Section~\ref{sec:projA}; H-A1; A-G1).
\item Calibrated uplift becomes the main rival for Project~B.
\item The gain is described as relative fit rather than formal evidence,
and the role of SLT in Project~A is stated as supporting.
\item Publication probabilities now include the chance that A-G1 fails.
The P0 note goes to the Alignment Forum first and to arXiv after the ICML
deadline, and the ICML rebuttal is scheduled.
\end{enumerate}

\paragraph{Revision 6.3 (23 September 2026): e-values, susceptibilities and new testbeds.}
\begin{enumerate}[leftmargin=*]
\item E-values are added as a statistical layer (Section~\ref{sec:evalues}).
For Project~A they separate a falsification regime, with prior-free
tests, from an Occam regime, where evidence is necessarily relative to a
prior (Section~\ref{sec:atests}; T5). For Project~B they give audits with
false-positive guarantees.
\item Susceptibilities are added (Section~\ref{sec:suscept}): a data-side
check for abstractions (H-A6), susceptibility clusters as candidate
variables, a prediction of elicitation cost from local samples (H-B5),
and a calibration suite for susceptibility estimates
(Section~\ref{sec:tools}).
\item Project~B gains a rare-output route using posterior proposals (H-B6),
funded from the reserve, and two organisms: sleeper-agent models and a
hidden-objective model.
\item A2 adds sparse-autoencoder and transcoder features as candidates.
Positioning adds ARC's surprise accounting and Resolution's susceptibility
and patterning work.
\item Effort, reserve rules, contacts and risks are updated accordingly.
\end{enumerate}

\paragraph{Revision 6.4 (23 September 2026): audit corrections.}
\begin{enumerate}[leftmargin=*]
\item The unlock-cost formula is stated consistently as
$k^*\approx D/(\beta\Delta\ell_C)$ in Section~\ref{sec:projB} and
Appendix~\ref{app:codes}, with $\delta$ as its local-sample estimate.
\item Claims are tightened where they ran ahead of the evidence. The
abstract presents the susceptibility and rare-output uses as tests; the
dilemma regime has a faithfulness gap close to, not exactly, zero;
universal inference requires a prequential numerator; the conformal audit
guarantee requires exchangeability with the benign library; and the
sleeper-agent and hidden-objective organisms are re-created on open models.
\item H-A6 moves next to H-A5, with a concrete observable, and a pointer to
H-B4 is added.
\item The positioning paragraph of Project~A is reordered; the notation
covers e-values, susceptibilities and $\delta$; the tooling row of
Section~\ref{sec:purpose} is updated; and the ICLR 2026 title of the
``Almost Bayesian'' paper is used.
\end{enumerate}

\section{Archive of completed checks}

The following scripts and outputs accompany this document. They concern
the stated mathematical models only.
\begin{description}[style=nextline,leftmargin=1.5em]
\item[\texttt{discrete\_sgd\_check.py}, \texttt{.json}] Monte Carlo of the
balanced discrete update: step-size sweep, label-noise case and fixed
$\eta/b$.
\item[\texttt{trapped\_check.py}] Non-exit fraction and final positions of
trapped paths under label noise.
\item[\texttt{tube\_mfpt.py}, \texttt{.json}] Finite-difference
mean-first-passage times in flat-set tubes.
\item[\texttt{identifiability.py}, \texttt{.json}] Collinearity of
logarithmic complexity terms against offsets.
\item[\texttt{thermometer\_check.py}, \texttt{.json}] Quadrature check of the loss-fluctuation thermometer.
\item[\texttt{batching\_check.py}, \texttt{.json}] Exact cumulative excess of Gaussian-mean learners under different block schemes.
\item[\texttt{topo\_interleave\_check.py}, \texttt{.json}] Interleaving bound
for bi-Lipschitz maps and the kernel identity for $\sep_V$.
\item[\texttt{topo\_bichromatic\_check.py}, \texttt{.json}] Closest
bichromatic pairs under geometry blind to $V$.
\item[\texttt{topo\_realizability\_check.py}, \texttt{.json}] Realizability
bound for counterfactual states and its diagram form.
\item[\texttt{topo\_circle\_check.py}, \texttt{.json}] Degree-one bars of
evenly spaced cyclic values. All four scripts import \texttt{common.py}.
\end{description}

\clearpage
\addcontentsline{toc}{section}{References}
\begin{thebibliography}{99}\small\raggedright
\setlength{\itemsep}{3pt}

\bibitem{adamaszek2017}\unv{} M.~Adamaszek and H.~Adams. The Vietoris--Rips
complexes of a circle. \emph{Pacific Journal of Mathematics} 290(1):1--40,
2017.

\bibitem{angell2026} R.~Angell, R.~Singhal, Z.~Horvitz, et al.
Estimating tail risks in language model output distributions. \emph{ICML
2026}; arXiv:2604.22167.

\bibitem{aoyagiwatanabe2005} M.~Aoyagi and S.~Watanabe. Stochastic
complexities of reduced rank regression in Bayesian estimation.
\emph{Neural Networks} 18(7):924--933, 2005.

\bibitem{arditi2024}\unv{} A.~Arditi, O.~Obeso, A.~Syed, D.~Paleka,
N.~Panickssery, W.~Gurnee and N.~Nanda. Refusal in language models is
mediated by a single direction. \emph{NeurIPS 2024}; arXiv:2406.11717.

\bibitem{ayonrinde2024}\unv{} K.~Ayonrinde, M.~T.~Pearce and L.~Sharkey.
Interpretability as compression: reconsidering SAE explanations of neural
activations with MDL-SAEs. Preprint, arXiv:2410.11179, 2024.

\bibitem{baker2026} G.~Baker, G.~Wang, J.~Hoogland, V.~Pathak and
D.~Murfet. Structural inference: interpreting small language models with
susceptibilities. \emph{ICLR 2026}; arXiv:2504.18274.

\bibitem{barannikov2022}\unv{} S.~Barannikov, I.~Trofimov, N.~Balabin and
E.~Burnaev. Representation topology divergence: a method for comparing neural
network representations. \emph{ICML 2022}; arXiv:2201.00058.

\bibitem{betley2026} J.~Betley, N.~Warncke, A.~Sztyber-Betley, D.~Tan,
X.~Bao, M.~Soto, M.~Srivastava, N.~Labenz and O.~Evans. Training large
language models on narrow tasks can lead to broad misalignment.
\emph{Nature} 649:584--589, 2026.

\bibitem{bhalla2026}\unv{} U.~Bhalla, E.~Bigelow, C.~Rager, T.~Haklay and
R.~Sarfati. Manifold steering reveals the shared geometry of neural network
representation and behavior. Preprint, arXiv:2605.05115, 2026.

\bibitem{bitsthatcount2026}\unv{} Bits that count: quantifying and
predicting capabilities of language models. ICML 2026 poster (authors not
verified).

\bibitem{blierollivier2018} L.~Blier and Y.~Ollivier. The description
length of deep learning models. \emph{NeurIPS 31}, 2018. arXiv:1802.07044.

\bibitem{brown2026} J.~R.~Brown, P.~Leask and L.~McKinney. Evil spectra:
how optimisers can amplify or suppress emergent misalignment. Preprint,
arXiv:2606.31591, 2026.

\bibitem{bushnaq2024}\unv{} L.~Bushnaq, J.~Mendel, et al. Using degeneracy
in the loss landscape for mechanistic interpretability. Preprint,
arXiv:2405.10927, 2024.

\bibitem{cai2026} Y.~Cai. Bypassing minimization bias: a shift-invariant
variance estimator for off-equilibrium local learning coefficients.
Preprint, arXiv:2606.22389, 2026.

\bibitem{chazal2014} F.~Chazal, V.~de Silva and S.~Oudot. Persistence
stability for geometric complexes. \emph{Geometriae Dedicata}
173:193--214, 2014. arXiv:1207.3885.

\bibitem{chazal2016} F.~Chazal, V.~de Silva, M.~Glisse and S.~Oudot.
\emph{The Structure and Stability of Persistence Modules}. SpringerBriefs in
Mathematics, Springer, 2016.

\bibitem{clarkebarron1990} B.~S.~Clarke and A.~R.~Barron.
Information-theoretic asymptotics of Bayes methods. \emph{IEEE
Transactions on Information Theory} 36(3):453--471, 1990.

\bibitem{cullen2026} B.~Cullen, S.~Estan-Ruiz, R.~Danait and J.~Li. A
basin-selection perspective on grokking via singular learning theory.
Preprint, arXiv:2603.01192 (v3, May 2026).

\bibitem{cultrera2025chromatic} S.~Cultrera di Montesano, O.~Draganov,
H.~Edelsbrunner and M.~Saghafian. Chromatic alpha complexes.
\emph{Foundations of Data Science} 8:30--62, 2025; arXiv:2212.03128.

\bibitem{dawid1984} A.~P.~Dawid. Present position and potential
developments: some personal views. Statistical theory: the prequential
approach. \emph{Journal of the Royal Statistical Society A}
147(2):278--292, 1984.

\bibitem{dinh2015}\unv{} L.~Dinh, D.~Krueger and Y.~Bengio. NICE: non-linear
independent components estimation. ICLR 2015 workshop; arXiv:1410.8516.

\bibitem{donoway2026} E.~Donoway, H.~Joren, F.~Roger and J.~Leike. Excess
description length of learning generalizable predictors. Preprint,
arXiv:2601.04728, 2026.

\bibitem{elicitation2025}\unv{} Quantifying elicitation of latent
capabilities in language models. NeurIPS 2025 poster (authors not
verified).

\bibitem{engels2025} J.~Engels, E.~J.~Michaud, I.~Liao, W.~Gurnee and
M.~Tegmark. Not all language model features are one-dimensionally linear.
\emph{ICLR 2025}; arXiv:2405.14860.

\bibitem{gauderis2026} W.~Gauderis, T.~Dooms, S.~T.~Holmer, K.~Ayonrinde
and G.~A.~Wiggins. From mechanistic to compositional interpretability.
ICML 2026; arXiv:2605.08934. (The third author's surname is spelled
differently across listings.)

\bibitem{geiger2024das}\unv{} A.~Geiger, Z.~Wu, C.~Potts, T.~Icard and
N.~D.~Goodman. Finding alignments between interpretable causal variables
and distributed neural representations. \emph{CLeaR 2024};
arXiv:2303.02536.

\bibitem{geiger2025ca}\unv{} A.~Geiger, D.~Ibeling, A.~Zur, et al. Causal
abstraction: a theoretical foundation for mechanistic interpretability.
Preprint, arXiv:2301.04709.

\bibitem{goldowskydill2025}\unv{} N.~Goldowsky-Dill, B.~Chughtai,
S.~Heimersheim and M.~Hobbhahn. Detecting strategic deception using
linear probes. \emph{ICML 2025}; arXiv:2502.03407.

\bibitem{gomez2017}\unv{} A.~N.~Gomez, M.~Ren, R.~Urtasun and R.~B.~Grosse.
The reversible residual network: backpropagation without storing
activations. \emph{NeurIPS 2017}; arXiv:1707.04585.

\bibitem{gordon2026} A.~Gordon, G.~Baker, G.~Wang, W.~Snell, S.~van
Wingerden and D.~Murfet. Towards spectroscopy: susceptibility clusters in
language models. Preprint, arXiv:2601.12703, 2026.

\bibitem{gozzi2025} F.~J.~Gozzi, M.~A.~Cerdeiro and P.~E.~Riera. Persistent
cost of Lipschitz maps. Preprint, arXiv:2511.07674 (v2, March 2026).

\bibitem{greenblatt2024}\unv{} R.~Greenblatt, F.~Roger, D.~Krasheninnikov
and D.~Krueger. Stress-testing capability elicitation with password-locked
models. \emph{NeurIPS 2024}; arXiv:2405.19550.

\bibitem{gross2024}\unv{} J.~Gross et al. Compact proofs of model
performance via mechanistic interpretability. \emph{NeurIPS 2024};
arXiv:2406.11779.

\bibitem{grunwald2007} P.~D.~Gr\"unwald. \emph{The Minimum Description
Length Principle}. MIT Press, 2007.

\bibitem{grunwald2024safe} P.~Gr\"unwald, R.~de Heide and W.~M.~Koolen.
Safe testing. \emph{Journal of the Royal Statistical Society B}
86(5):1091--1128, 2024.

\bibitem{hennick2026} M.~Hennick and S.~De Baerdemacker. Almost Bayesian: dynamics of SGD through singular learning theory. \emph{ICLR 2026};
arXiv:2503.22478 (v2, March 2026).

\bibitem{hewittliang2019} J.~Hewitt and P.~Liang. Designing and
interpreting probes with control tasks. \emph{EMNLP 2019};
arXiv:1909.03368.

\bibitem{hilton2024}\unv{} J.~Hilton. Formal verification, heuristic
explanations and surprise accounting. Alignment Research Center blog,
2024.

\bibitem{hirsch1976} M.~W.~Hirsch. \emph{Differential Topology}. Graduate
Texts in Mathematics 33. Springer, 1976.

\bibitem{hitchcock2025} R.~Hitchcock and J.~Hoogland. From global to local:
a scalable benchmark for local posterior sampling. Preprint,
arXiv:2507.21449 (v2, August 2026).

\bibitem{hubinger2024}\unv{} E.~Hubinger et al. Sleeper agents: training
deceptive LLMs that persist through safety training. Preprint,
arXiv:2401.05566, 2024.

\bibitem{interpbench2024} R.~Gupta, I.~Arcuschin, T.~Kwa and
A.~Garriga-Alonso. InterpBench: semi-synthetic transformers for evaluating
mechanistic interpretability techniques. \emph{NeurIPS 2024 Datasets and
Benchmarks}; arXiv:2407.14494.

\bibitem{jammalamadaka1986}\unv{} S.~R.~Jammalamadaka and S.~Janson. Limit
theorems for a triangular scheme of U-statistics with applications to
inter-point distances. \emph{The Annals of Probability} 14(4):1347--1358,
1986.

\bibitem{kreer2026} P.~A.~Kreer, W.~Wu, M.~Adam, Z.~Furman and J.~Hoogland.
Bayesian influence functions for Hessian-free data attribution.
\emph{ICLR 2026}; arXiv:2509.26544.

\bibitem{larsson2026} M.~Larsson. Langevin dynamics along the zero set of
real-analytic potentials. Preprint, arXiv:2608.09840, 2026.

\bibitem{lau2025llc} E.~Lau, Z.~Furman, G.~Wang, D.~Murfet and S.~Wei. The
local learning coefficient: a singularity-aware complexity measure.
\emph{AISTATS 2025}, PMLR 258; arXiv:2308.12108.

\bibitem{li2024wmdp}\unv{} N.~Li, A.~Pan, et al. The WMDP benchmark:
measuring and reducing malicious use with unlearning. \emph{ICML 2024};
arXiv:2403.03218.

\bibitem{lucki2024}\unv{} J.~\L{}ucki, B.~Wei, Y.~Huang, P.~Henderson,
F.~Tram\`er and J.~Rando. An adversarial perspective on machine
unlearning for AI safety. Preprint, arXiv:2409.18025, 2024.

\bibitem{ma2026} S.~Ma, Y.~Luo, Zhangji, C.~Xiao, A.~Gao, W.-H.~Huang, W.~Wang,
Q.~Wu, X.~Li, J.~Wei and Q.~Zhang. Hidden APIs in language models:
discovering reusable causal interfaces from forked futures. Preprint,
arXiv:2607.27617, 2026.

\bibitem{makelov2024}\unv{} A.~Makelov, G.~Lange and N.~Nanda. Is this the
subspace you are looking for? An interpretability illusion for subspace
activation patching. \emph{ICLR 2024}; arXiv:2311.17030.

\bibitem{mandt2017} S.~Mandt, M.~D.~Hoffman and D.~M.~Blei. Stochastic
gradient descent as approximate Bayesian inference. \emph{Journal of
Machine Learning Research} 18(134):1--35, 2017.

\bibitem{marks2025}\unv{} S.~Marks et al. Auditing language models for
hidden objectives. Preprint, arXiv:2503.10965, 2025.

\bibitem{meloux2025ident}\unv{} M.~M\'eloux, S.~Maniu, F.~Portet and
M.~Peyrard. Everything, everywhere, all at once: is mechanistic
interpretability identifiable? \emph{ICLR 2025}; arXiv:2502.20914.

\bibitem{meloux2026} M.~M\'eloux, T.~Pimentel, F.~Portet and M.~Peyrard.
Validating causal abstraction metrics on simulated complex systems.
Preprint, arXiv:2607.00267, 2026.

\bibitem{mib2025} A.~Mueller, A.~Geiger, S.~Wiegreffe, et al. MIB: a
mechanistic interpretability benchmark. \emph{ICML 2025};
arXiv:2504.13151.

\bibitem{naitzat2020} G.~Naitzat, A.~Zhitnikov and L.-H.~Lim. Topology of
deep neural networks. \emph{Journal of Machine Learning Research}
21(184):1--40, 2020.

\bibitem{patterningpractice2026}\unv{} Patterning in practice: debiasing
reward models with susceptibilities. Preprint, arXiv:2609.00699, 2026
(authors not verified).

\bibitem{phalign2026}\unv{} Tracking representation dynamics in large
language models with persistent homology. Preprint, arXiv:2606.19542, 2026
(authors not verified).

\bibitem{pimentel2020info} T.~Pimentel, J.~Valvoda, R.~H.~Maudslay,
R.~Zmigrod, A.~Williams and R.~Cotterell. Information-theoretic probing
for linguistic structure. \emph{ACL 2020}; arXiv:2004.03061.

\bibitem{pimentel2020pareto} T.~Pimentel, N.~Saphra, A.~Williams and
R.~Cotterell. Pareto probing: trading off accuracy for complexity.
\emph{EMNLP 2020}; arXiv:2010.02180.

\bibitem{ramdas2023}\unv{} A.~Ramdas, P.~Gr\"unwald, V.~Vovk and G.~Shafer.
Game-theoretic statistics and safe anytime-valid inference.
\emph{Statistical Science} 38(4):576--601, 2023.

\bibitem{reguera2001} D.~Reguera and J.~M.~Rub\'i. Kinetic equations for
diffusion in the presence of entropic barriers. \emph{Physical Review E}
64:061106, 2001.

\bibitem{rissanen1986} J.~Rissanen. Stochastic complexity and modeling.
\emph{The Annals of Statistics} 14(3):1080--1100, 1986.

\bibitem{shtarkov1987} Yu.~M.~Shtarkov. Universal sequential coding of
single messages. \emph{Problems of Information Transmission}
23(3):175--186, 1987.

\bibitem{soligo2026} A.~Soligo, E.~Turner, S.~Rajamanoharan and N.~Nanda.
Emergent misalignment is easy, narrow misalignment is hard. \emph{ICLR
2026}; arXiv:2602.07852.

\bibitem{sutter2025} D.~Sutter, J.~Minder, T.~Hofmann and T.~Pimentel. The
non-linear representation dilemma: is causal abstraction enough for
mechanistic interpretability? \emph{NeurIPS 2025}; arXiv:2507.08802.

\bibitem{taylor2025} J.~Taylor, S.~Black, D.~Bowen, et al. Auditing games
for sandbagging. Preprint, arXiv:2512.07810, 2025.

\bibitem{teshima2020}\unv{} T.~Teshima, I.~Ishikawa, K.~Tojo, K.~Oono,
M.~Ikeda and M.~Sugiyama. Coupling-based invertible neural networks are
universal diffeomorphism approximators. \emph{NeurIPS 2020};
arXiv:2006.11469.

\bibitem{tice2025} C.~Tice, P.~A.~Kreer, N.~Helm-Burger, P.~S.~Shahani,
F.~Ryzhenkov, J.~Haimes, F.~Hofst\"atter and T.~van der Weij. Noise
injection reveals hidden capabilities of sandbagging language models.
\emph{NeurIPS 2025}; arXiv:2412.01784.

\bibitem{tracr2023}\unv{} D.~Lindner, J.~Kram\'ar, S.~Farquhar, M.~Rahtz,
T.~McGrath and V.~Mikulik. Tracr: compiled transformers as a laboratory
for interpretability. \emph{NeurIPS 2023}; arXiv:2301.05062.

\bibitem{turner2025}\unv{} E.~Turner, A.~Soligo, et al. Model organisms
for emergent misalignment. Preprint, arXiv:2506.11613, 2025.

\bibitem{urdshals2025} E.~Urdshals, E.~Lau, J.~Hoogland, S.~van Wingerden
and D.~Murfet. Compressibility measures complexity: minimum description
length meets singular learning theory. Preprint, arXiv:2510.12077, 2025.

\bibitem{vanderweij2024}\unv{} T.~van der Weij, F.~Hofst\"atter, et al. AI
sandbagging: language models can strategically underperform on
evaluations. \emph{ICLR 2025}; arXiv:2406.07358.

\bibitem{vershynin2018} R.~Vershynin. \emph{High-Dimensional Probability}.
Cambridge University Press, 2018.

\bibitem{voita2020} E.~Voita and I.~Titov. Information-theoretic probing
with minimum description length. \emph{EMNLP 2020}; arXiv:2003.12298.

\bibitem{vovk2005} V.~Vovk, A.~Gammerman and G.~Shafer. \emph{Algorithmic
Learning in a Random World}. Springer, 2005.

\bibitem{wangmurfet2026} G.~Wang and D.~Murfet. Patterning: the dual of
interpretability. \emph{ICML 2026}; arXiv:2601.13548.

\bibitem{wasserman2020}\unv{} L.~Wasserman, A.~Ramdas and S.~Balakrishnan.
Universal inference. \emph{Proceedings of the National Academy of
Sciences} 117(29):16880--16890, 2020.

\bibitem{watanabe2009} S.~Watanabe. \emph{Algebraic Geometry and
Statistical Learning Theory}. Cambridge University Press, 2009.

\bibitem{watanabe2013} S.~Watanabe. A widely applicable Bayesian information
criterion. \emph{Journal of Machine Learning Research} 14(27):867--897,
2013.

\bibitem{watanabe2018} S.~Watanabe. \emph{Mathematical Theory of Bayesian
Statistics}. CRC Press, 2018.

\bibitem{wuhilton2025} G.~Wu and J.~Hilton. Estimating the probabilities
of rare outputs in language models. \emph{ICLR 2025}; arXiv:2410.13211.

\bibitem{xu2020} Y.~Xu, S.~Zhao, J.~Song, R.~Stewart and S.~Ermon. A theory
of usable information under computational constraints. \emph{ICLR 2020};
arXiv:2002.10689.

\bibitem{zwanzig1992} R.~Zwanzig. Diffusion past an entropy barrier.
\emph{The Journal of Physical Chemistry} 96(10):3926--3930, 1992.

\end{thebibliography}
\end{document}
